\documentclass[11pt]{article}
\usepackage[margin=1.1in]{geometry}
\usepackage{amsmath,amssymb,amsthm}
\usepackage{booktabs}
\usepackage[colorlinks=true,linkcolor=blue,citecolor=blue]{hyperref}
\numberwithin{equation}{section}
\newtheorem{theorem}{Theorem}[section]
\newtheorem{lemma}[theorem]{Lemma}
\newtheorem{proposition}[theorem]{Proposition}
\newtheorem{corollary}[theorem]{Corollary}
\theoremstyle{definition}
\newtheorem{hypothesis}[theorem]{Hypothesis}
\theoremstyle{remark}
\newtheorem{remark}[theorem]{Remark}
\newcommand{\LL}{\Lambda}
\newcommand{\SSS}{\mathfrak{S}}
\newcommand{\e}{\mathrm{e}}
\newcommand{\Si}{\mathrm{Si}}
\title{\bf More than $0.7947$ of the zeros of the Riemann zeta
function are simple and on the critical
line\thanks{Verification status (read this first): the chain of
this paper consists of classical inputs, machine-verified
identities, certified constants, and numerically evaluated
connected constants carrying stated two-sided bands. Its claim
grade is \emph{certified-candidate}, not record: the certified
closed forms of $C_5$, $\{4,2\}$ and $\{6\}$ remain open ---
though each is proved RATIONAL (a finite signed sum of rational
polytope volumes, \S\ref{s:conv}), giving the designated
exact-integration route an explicit representation --- the Lean
formalization covers the identity and certificate layers only
(kernel-compiled, core Lean~4, zero \texttt{sorry};
\S\ref{s:verif}(f)), the analytic layer is unformalized, and no
external review has taken place. The three write-out items
W1--W3 are written out (\S\ref{s:writeouts}), the pairing-layer
constants are certified as exact rationals (\S\ref{s:five:cert},
Appendix~\ref{a:cert}), and the consumption step is backed by an
exact rational dual certificate valid on unbounded support
(Lemma~\ref{l:dual}). The precise grading of every layer, the
discrepancy register, and the program's verification gate
(formalization and external review precede any record claim) are
in \S\ref{s:verif} and Appendix~\ref{a:record}.}}
\author{Hongyi Yang\\[2pt]
\small University of Toronto\\
\small \texttt{hongy.yang@mail.utoronto.ca}
\and Shihua Yang\\[2pt]
\small The University of Hong Kong\\
\small \texttt{u3590485@connect.hku.hk}}
\date{August 2026}
\begin{document}
\maketitle

\begin{abstract}
Let $N(T,2T)$ count the nontrivial zeros of $\zeta(s)$ with
ordinates in $(T,2T]$, and $N_0^{s}$, $N_d$ those simple and on
the critical line, respectively distinct. Subject to the
verification grading stated below and in \S\ref{s:verif}, we
establish, with ineffective constants,
\[
\liminf_{T\to\infty}\frac{N_0^{s}}{N}\;\ge\;0.7947,
\qquad
\liminf_{T\to\infty}\frac{N_d}{N}\;\ge\;0.8973,
\]
and the same for any fixed primitive Dirichlet $L$-function; en
route, the rungs $0.6916/0.8458$ (one-sided fourth-moment
remainder, via Matom\"aki--Radziwi\l{}\l--Tao),
and $13/18=0.7222$, $31/36=0.8611$ (exact fourth moment). The proof
evaluates
the third through sixth trace moments of the Gabor-compressed
Weil matrix of \cite{C26} --- $m_3=2$ and $m_4=\tfrac{13}4$
exactly, $m_5=\tfrac{67}{12}+C_5$ and $m_6=10.1657(18)$, the
remainders being connected sine-kernel Ursell constants --- and consumes them through a
moment-pinned linear program built on an explicit
Chebyshev--Markov certificate --- an
\emph{exact rational} dual certificate, valid on unbounded
support. The fourth-moment engine (``Lemma D'') is proved in the
truncated form the program consumes: the cell ledger is
$\ell_1$-truncated at $(\log X)^{B}$ with power-saving tails, the
lock variance transports by exact product and Parseval identities
to a spectral frame in which every factor is a windowed single
prime sum at a rational of logarithmic denominator, and the
Siegel--Walfisz theorem plus Vaughan's estimate close all bins
--- no dispersion method and no large-modulus technology
anywhere. The local main terms at every cycle length are pinned
by a four-line coincidence-counting closed form; the Bell$(5)$
and Bell$(6)$ ledgers collapse onto four-cycle anchors, now
certified as exact rationals ($t_{\mathrm{adj}}=\tfrac7{60}$,
$t_{\mathrm{opp}}=\tfrac1{30}$,
$\{2,2,2\}=\tfrac{131}{420}$: exact polyhedral integration,
Appendix~\ref{a:cert}), through machine-exact frozen-slot
identities. \emph{Claim grading}: the identity, certificate and
pairing-constant layers are machine-verified or exact-rational;
the write-out items W1--W3
are written out in \S\ref{s:writeouts}; the analytic chain is
certified-candidate; the three connected
constants $C_5=0.0278(1)$, $\{4,2\}=-0.0552(8)$ and
$\{6\}=-0.0078(8)$ are midpoint-grid ladder evaluations with
stated bands (\S\ref{s:conv}), proved rational by the polytope
dimension theorem of \S\ref{s:conv} but not yet evaluated in
closed form; the consumption layer (inertia, tail, counting)
is that of \cite{C26}, recalled in \S\ref{s:frame}.
The consumption step itself is exact: a rational sum-of-squares
dual certificate, kernel-checked in core Lean~4 (twenty theorems across the
identity, local-law and certificate layers, zero
\texttt{sorry}). A standalone reproduction and
certification package accompanies the paper
(\url{https://github.com/JoshuaHKU/zeta-0.7947-reproduction});
a falsifier registry
($194$ pre-registered checks, $34$ fired and converted, all
forward-recorded) documents the research trail.
\end{abstract}

\tableofcontents

\section{Introduction}\label{s:intro}

That a positive proportion of the nontrivial zeros of $\zeta(s)$
lie on the critical line is due to Selberg \cite{Sel42}; Levinson
\cite{Lev74} obtained the proportion $\tfrac13$, Conrey
\cite{Con89} reached $\tfrac25$, and the record along Levinson's
method stands at $\kappa>5/12$ \cite{PRZZ} (see also
\cite{BCY11,Fen12}). Very recently the two-thirds theorem
\cite{C26} proceeded by a different route: a linear-algebraic
reading of Montgomery's pair-correlation sum, in which spectral
information is read from trace moments of a Gabor compression
$\widetilde G$ of Weil's Hermitian form, and the first two
moments alone --- both computable unconditionally from
Montgomery's prime-side evaluation at bandwidth $\lambda\le1$
\cite{Mon73,BGST} --- already yield the proportion $2/3-o(1)$
(and $0.6725$ with the Montgomery--Taylor window \cite{Mon75}).
The present paper extends that framework to the third through
sixth moments. By \cite[Rem.~1.1]{C26} no certificate consuming
only the bandwidth-one data can exceed $0.68185$; the route past
the ceiling identified there (\S7.5(f) of \cite{C26}) is the
fourth moment,
$\operatorname{tr}\widetilde G^{4}=d\ell_1^{4}(13/4+o(1))$, worth
$13/18$.

This paper climbs that ladder in four stages, in a single
text with one notation, one bibliography, one verification
ledger, and one claim grading: the certificate calculus and the
third moment (\S\ref{s:cert}), a logically independent
one-sided fourth-moment route reaching $0.6916$ (\S\ref{s:fm}),
the exact fourth moment worth $13/18$
(\S\S\ref{s:trunc}--\ref{s:four}), and the fifth and sixth
moments with their consumption
(\S\S\ref{s:five}--\ref{s:lp}), reaching the theorem
below.

\begin{theorem}[Main theorem; grade: certified-candidate, see
\S\ref{s:verif}]\label{t:main}
With ineffective constants,
\[
\liminf_{T\to\infty}\frac{N_0^{s}(T,2T)}{N(T,2T)}\ge 0.7947,
\qquad
\liminf_{T\to\infty}\frac{N_d(T,2T)}{N(T,2T)}\ge 0.8973,
\]
and the same holds for the zeros of any fixed primitive Dirichlet
$L$-function. En route, each rung consuming exactly what its
chain establishes: $m_4\to13/4$ two-sidedly (hence
$13/18-\varepsilon$ and $31/36-\varepsilon$), and, by the
logically independent one-sided route of \S\ref{s:fm},
$0.6916/0.8458$. (An intermediate fourth-and-fifth-moment rung considered at
an earlier stage of this research is repaired and demoted in
\S\ref{s:lp}: without sixth-moment information it is not
supported on unbounded spectral support --- register item D6.)
\end{theorem}

The dependency structure of Theorem~\ref{t:main} is graded
explicitly in \S\ref{s:verif}: the $13/18$ rung rests on the
classical-technology proof of Lemma D (\S\ref{s:four}) and
the consumption layer of \cite{C26} (recalled in
\S\ref{s:frame}); the $0.7947$ headline
additionally consumes the three connected Ursell constants of
\S\ref{s:five} as numerical evaluations with stated two-sided
bands (the pairing layer is exact, \S\ref{s:five:cert}). Nothing in this paper bears on the Riemann Hypothesis
itself: proportion statements are insensitive to $o(N)$ off-line
zeros: the inputs (functional equation, explicit formula, mean
values of Dirichlet polynomials of length $\le T$) are satisfied
by objects for which the analogue of RH is false, and the
argument produces lower bounds only \cite[\S1.5]{C26}.

\subsection*{Relation to \cite{C26}, and what is new}

This paper is built on \cite{C26} and consumes its framework
without repetition: the compressed Weil matrix, the zero-side
block structure, the tail estimate, and the first two trace
moments are used exactly as established there
(\S\ref{s:frame}). The advance is fourfold. \emph{(i) The
moment ladder}: \cite{C26} consumes only
$\mu_1\to1$, $\mu_2\to\tfrac43$; this paper evaluates
$\mu_3\to2$ and $\mu_4\to\tfrac{13}4$ exactly, and
$\mu_5$, $\mu_6$ up to three connected constants with proved
rationality (\S\S\ref{s:cert}--\ref{s:five}). \emph{(ii)
The truncation technology}: the $\ell_1$-truncated cell ledger
with Siegel--Walfisz/Vaughan closure (\S\ref{s:trunc}), absent
from \cite{C26}, is what makes the higher moments computable
with classical tools. \emph{(iii) The consumption}: the
rank--trace inequality of \cite[Lem.~3.2]{C26} --- the matrix
form of $m^{2}\ge2m-1$ --- is replaced by an exact
Chebyshev--Markov dual certificate at degree six
(\S\ref{s:lp}), the analogue of $m^{2}\ge3m-2$ and beyond,
valid on unbounded support with exact rational corners.
\emph{(iv) The verification apparatus}: pre-registered
falsifier checks, a forward-only discrepancy register, and a
kernel-compiled Lean layer (\S\ref{s:verif}), replacing the
traditional single-grade presentation of \cite{C26}. The
constants rise from $2/3$ and $\tfrac56$ (distinct) to
$0.7947$ and $0.8973$.

\subsection*{Relation to the pair-correlation route}

Baluyot, Goldston, Suriajaya and Turnage-Butterbaugh \cite{BGST}
showed that Montgomery's pair-correlation method yields
two-thirds-type conclusions under hypotheses weaker than RH, and
Goldston--Suriajaya \cite{GS} reach the two-thirds constant
itself conditionally; those are conditional statements about the
\emph{zeros'} pair correlation. The present route is disjoint,
and unconditional: it consumes the average behaviour of
$2k$-point correlations of $\LL\star\LL$ --- prime-side
objects, controlled by the Siegel--Walfisz theorem and Vaughan's
estimate alone --- through the compressed-matrix framework, with
no hypothesis on the zeros. The two-thirds threshold, previously
reached only under such hypotheses, is crossed here
unconditionally with margin ($0.7947>2/3$), and the additional
information consumed relative to the bandwidth-one data is
exactly the third-through-sixth moment chain of
\S\S\ref{s:cert}--\ref{s:five}.

\subsection*{The method in one page}

The moments $m_k$ of the spectral measure of $\widetilde G$ are
carried by $k$-cycle ledgers: cells indexed by prime-power
modulus tuples and lock (offset) vectors, weighted
$1/(d\ell_1^{k})$, with window-overlap volumes supplying a walk
geometry. The consumption is an explicit Chebyshev--Markov
certificate (\S\ref{s:cert}): mass near the origin of the
spectral measure is bounded by moment data alone, with exact
exchange rates. Three observations reduce the moment evaluations
to classical tools.

\emph{Truncation (\S\ref{s:trunc}).} The $1/\ell_1^{k}$ weights
discharge all cells with $\ell_1>(\log X)^{B}$ trivially and
two-sidedly; only logarithmic moduli remain.

\emph{The spectral frame (\S\ref{s:four}, \S\ref{s:five}).} The
lock variance of a $k$-cycle obeys exact identities --- at $k=4$
a discrete Parseval identity; at every $k$ the \emph{product
identity} $\sum_j N_k(j)^2=\sum_t A_1(t)\cdots A_k(t)$ (two
chains share a lock vector iff they differ by a common shift) ---
which exhibit it as a spectral $L^2$-distance whose every factor
is a \emph{windowed single prime sum}. At logarithmic moduli
every main-term frequency is a rational of logarithmic
denominator: the Siegel--Walfisz theorem gives an arbitrary
log-power saving pointwise on the doubly-major bins, Vaughan's
estimate closes the minor and mixed bins, and the singular-series
comb is identified exactly by the local law below. The apparent
``twin hardness'' of locked correlations dissolves: no per-lock
evaluation is ever attempted.

\emph{The local law and the anchor collapse (\S\ref{s:local},
\S\ref{s:five}).} The mod-$p$ structure of every cycle length is
an affine lock chain, whose density obeys a coincidence-counting
closed form proved in four lines and verified exhaustively. The
Bell$(5)=52$ and Bell$(6)=203$ class ledgers of the fifth and
sixth moments then collapse onto the certified four-cycle
anchors through machine-exact frozen-slot identities; the only
genuinely new constants above $m_4$ are the connected Ursell
values.

\subsection*{Verification discipline and contributions}

The research program behind this paper operates falsifier-first:
every structural claim carries a pre-registered numerical check;
corrections are recorded forward, never edited away; claims are
graded. Thirty-four of the registry's checks have fired and been
converted --- several of them shaping this paper directly
(Appendix~\ref{a:record}). The chain's grade, its named
write-out items, and the results of the certification audit are
stated in \S\ref{s:verif}; per the program's verification gate,
Lean formalization and external review precede any record claim.
The mathematical development of this paper was carried out by
Claude (Anthropic), a large language model, within a research
program directed by the authors, who specified objectives,
verification discipline, and publication decisions, and who take
full academic responsibility for the content. In accordance with
prevailing journal authorship policies, the model is not listed
as an author; its role is stated here and in the
Acknowledgements. The theorem claims of this paper are advanced
\emph{at the stated grade}: as certified candidates, not as
established records.

\section{Framework and notation, following
\cite{C26}}\label{s:frame}

The construction of the compressed Weil matrix $\widetilde G$
--- Weil's Hermitian form, the test family, the zero-side block
structure, the tail estimate, and the first two trace moments
--- is exactly that of the two-thirds paper \cite{C26}
(\S\S2--5 there). This section recalls the notation and states
the five results of \cite{C26} that the moment ladder consumes;
proofs are in \cite{C26}, whose own analytic imports are the
published literature \cite{Wei52,IK,Tit86,Mon73,BGST}. Our
contribution begins with the third moment
(\S\ref{s:cert}).

\subsection{Weil's form, the test family, and
$\widetilde G$}\label{s:frame:G}

Throughout, $l=\log(T/2\pi)$, $\ell_1=l+2\log2-1$ (so
$N(T,2T)=T\ell_1/2\pi+O(l)$), and for fixed
$\lambda\in(0,1]$: $L=\lambda l$, $X=e^{L}$,
$d=\lfloor LT/2\pi\rfloor$. For $f\in C_c^{2}(\mathbb R)$
put $h_f(z)=\widehat f(z)=\int f(u)e^{izu}\,du$; two
integrations by parts give
\begin{equation}\label{e:hdecay}
|h_f(x+iy)|\le e^{|y|\Lambda_f}
\min\bigl(\|f\|_1,\|f''\|_1|x+iy|^{-2}\bigr),
\qquad\operatorname{supp}f\subset[-\Lambda_f,\Lambda_f].
\end{equation}
Weil's Hermitian form is
$W(f,g)=\sum_{\rho}m_\rho h_f(\gamma_\rho)
\overline{h_g(\gamma_\rho)}$, the sum over distinct
nontrivial zeros $\rho=\beta+i\gamma_\rho$ (multiplicity
$m_\rho$ explicit; absolutely convergent by \eqref{e:hdecay}
and $\sum m_\rho|\gamma_\rho|^{-2}<\infty$). The zero
multiset is invariant under $\rho\mapsto1-\bar\rho$, so $W$
is Hermitian. With $s=\tfrac12+i\tau$ define the prime-side
densities
\[
\mu(\tau)=\frac1{2\pi}\mathrm{Re}\,
\frac{\Gamma'}{\Gamma}\Bigl(\frac14+\frac{i\tau}2\Bigr)
-\frac{\log\pi}{2\pi},\quad
\Pi_X(\tau)=\frac{1}{2\pi(\tfrac14+\tau^{2})}
+\frac1\pi\,\mathrm{Re}\,\frac{X^{s}-1}{s},\quad
P_X(\tau)=-\frac1\pi\sum_{n\le X}
\frac{\LL(n)}{\sqrt n}\cos(\tau\log n),
\]
and $\nu_X=\mu+\Pi_X+P_X$. Weil's explicit formula
\cite{Wei52} (in the normalization of
\cite[Thm.~5.12]{IK}, applied to $k=f*\tilde g$,
$\tilde g(u)=g(-u)$, whose transform $h_k=h_fh_{\bar g}$ is
entire with the decay \eqref{e:hdecay}) gives the spectral
form
\begin{equation}\label{e:weil}
W(f,g)=\int_{\mathbb R}h_f(\tau)\overline{h_g(\tau)}\,
\nu_X(\tau)\,d\tau,
\qquad\operatorname{supp}f,g\subset[-\tfrac L2,\tfrac L2].
\end{equation}
We record $|\Pi_X(\tau)|\le3\sqrt X/(1+|\tau|)$,
$|P_X(\tau)|\ll\sqrt X$, $\int_T^{2T}\mu=N(T,2T)+O(l)$ and
$\int_T^{2T}\mu^{2}=T\ell_1^{2}/4\pi^{2}+O(l^{2})$
(Stirling).

Fix a nondecreasing $\varrho\in C^{3}(\mathbb R)$ with
$\varrho=0$ on $(-\infty,0]$, $\varrho=1$ on $[1,\infty)$,
a ramp width $1\le w\le L/8$, and the taper
$\phi(u)=\varrho\bigl((L/2-|u|)/w\bigr)$: even,
$0\le\phi\le1$, $\operatorname{supp}\phi=[-L/2,L/2]$,
$\phi=1$ on $[-L/2+w,L/2-w]$. Put $a=\tfrac1L\int\phi^{2}$,
$b=\tfrac1L\int\phi^{4}$ (so $1-2w/L\le b\le a\le1$),
$\Phi=\widehat{\phi^{2}}$, and note the window calculus
\begin{equation}\label{e:windowcalc}
(L-2w-|y|)_+\;\le\;g(y):=(\phi^{2}\star\phi^{2})(y)
\;\le\;(L-|y|)_+ .
\end{equation}
The Gabor system at critical density is
$f_k(u)=\phi(u)e^{-i\tau_ku}$, $\tau_k=T+2\pi k/L$,
$0\le k<d$, and the compressed Weil matrix is
\[
G_{kl}=W(f_k,f_l)
=\sum_\rho m_\rho\,\widehat\phi(\gamma_\rho-\tau_k)\,
\widehat\phi(\gamma_\rho-\tau_l)
=\int_{\mathbb R}\widehat\phi(\tau-\tau_k)
\widehat\phi(\tau-\tau_l)\,\nu_X(\tau)\,d\tau,
\qquad\widetilde G=G/L,
\]
real symmetric by either expression; its normalized trace
moments are
$\mu_k=\operatorname{tr}\widetilde G^{k}/(d\ell_1^{k})$.

\begin{lemma}[Poisson summation for the Gabor
system]\label{l:gabor}
For all $\tau,\tau'\in\mathbb R$,
$\sum_{k\in\mathbb Z}\widehat\phi(\tau-\tau_k)
\widehat\phi(\tau'-\tau_k)=L\,\Phi(\tau-\tau')$; in
particular $\sum_{k\in\mathbb Z}
\widehat\phi(\tau-\tau_k)^{2}=aL^{2}$.
\end{lemma}

\begin{proof}
This is \cite[Lem.~2.2]{C26}: the Gabor system at critical
density $h=2\pi/L$ has a translation-invariant frame kernel
with no aliasing error, for any window supported in an interval
of length $L$.
\end{proof}

\subsection{The zero side: inertia, tail, and the counting
inequality}\label{s:frame:count}

Fix $D_0=T^{1/2}$, $I=(T,2T]$, $I'=(T-D_0,2T+D_0]$, and split
$G=A+E$ over distinct zeros with $\gamma\in I'$ (matrix $A$)
and the rest ($E$); both are real symmetric. Classify the
distinct zeros with $\gamma\in I'$: $S_1$ (simple, on-line;
$s_1$), $S_2$ (multiple, on-line; $s_2$), and off-line pairs
$\{\rho,1-\bar\rho\}$ ($p$ of them), so
$\#Z(I')=s_1+s_2+2p$ and $N(I')\ge s_1+2s_2+2p$.

\begin{proposition}[Block structure]\label{p:inertiaA}
$n_+(\widetilde A)\le s_1+s_2+p$ and
$\operatorname{rank}\widetilde A\le\#Z(I')$.
\end{proposition}

\begin{proof}
This is \cite[Prop.~4.1]{C26}: Sylvester inertia under
pull-back --- an on-line zero contributes a positive $1\times1$
block, an off-line pair a hyperbolic block of signature $(1,1)$,
and no independence of the evaluation functionals is needed.
\end{proof}

\begin{proposition}[Tail]\label{p:tail}
Let $A_0$ be an absolute constant with
$N(t+1)-N(t)\le A_0\log(t+3)$ \cite[Thm.~9.2]{Tit86}. Then
for $T\ge T_0$,
\[
\|\widetilde E\|_{\mathrm{op}}\;\le\;\theta_0
:=4A_0C_1^{2}X^{1/2}\,\frac{\log(4T)}{D_0^{2}}
\;\ll\;l\,T^{\lambda/2-1},
\qquad C_1=\|\phi''\|_1=\frac{2\|\varrho''\|_1}{w},
\]
and the trace norm satisfies
$\|\widetilde E\|_1\le2\theta_0/L$.
\end{proposition}

\begin{proof}
This is \cite[Prop.~4.2]{C26} (the taper's $r^{-2}$-decay
\eqref{e:hdecay} against the zero-density bound); we restate the
explicit $\theta_0$ because the truncation discharge of
\S\ref{s:trunc} and the third-moment tail step consume it.
\end{proof}

\begin{proposition}[Counting inequality]\label{p:count}
For every $\theta\ge\theta_0$,
\[
s_1\;\ge\;2\,n_+^{\theta}(\widetilde G)-N(I'),
\qquad
\#Z(I')\;\ge\;n_+^{\theta}(\widetilde G),
\]
hence, with $N(I'\setminus I)\ll D_0l=o(N)$,
\begin{equation}\label{e:count}
N_0^{s}(T,2T)\ \ge\ 2\,n_+^{\theta}(\widetilde G)
-N(T,2T)-2N(I'\setminus I),
\qquad
N_d(T,2T)\ \ge\ n_+^{\theta}(\widetilde G)
-N(I'\setminus I).
\end{equation}
\end{proposition}

\begin{proof}
This is \cite[Prop.~4.5]{C26}: Weyl's inequality (at threshold
$\theta\ge\theta_0$, Proposition~\ref{p:tail}) reduces
$n_+^{\theta}(\widetilde G)$ to $n_+(\widetilde A)$, and
Proposition~\ref{p:inertiaA} with the multiplicity count
$s_2+p\le\tfrac12(N(I')-s_1)$ gives both displays.
\end{proof}

These inequalities hold for any locally finite multiset in the
strip invariant under $\rho\mapsto1-\bar\rho$ with
$N(t+1)-N(t)\ll\log(t+3)$; they contain no arithmetic. All
the arithmetic is in the trace moments.

\subsection{The first two moments, and
notation for the ladder}\label{s:frame:moments}

\begin{proposition}[First two moments]\label{p:mu12}
$\operatorname{tr}\widetilde G=aL\,N(T,2T)+O(L\sqrt X)$, and
\[
\operatorname{tr}\widetilde G^{2}
=2\pi bL\int_T^{2T}\mu^{2}
+\frac{T}{\pi}\sum_{n\le X}\frac{\LL(n)^{2}}{n}\,
g(\log n)+O\bigl(Ll\log l\,(l^{2}+X)\bigr).
\]
Consequently, in the endpoint regime $\lambda\to1^-$,
$w/L\to0$: $\mu_1\to1$ and $\mu_2\to\tfrac43$.
\end{proposition}

\begin{proof}
This is \cite[Prop.~5.3 and Thm.~5.8]{C26}: the trace by the
Gabor--Poisson collapse, the second moment by Montgomery's
prime-side evaluation of the pair-correlation second moment at
bandwidth $\le1$, unconditional in the band
\cite{Mon73,BGST}. We restate it because the diagonal
$\sum_{n\le X}(\LL(n)^{2}/n)\,g(\log n)
=\tfrac{L^{3}}{6}(1+O(1/L))$, with the same window
autocorrelation \eqref{e:windowcalc}, is the calibration
consumed by the third-moment ledger (Theorem~\ref{t:m3}) and
by the model gates of \S\ref{s:cert:model}.
\end{proof}

\begin{remark}[Endpoint conventions]\label{r:endpoint}
The endpoint $\lambda=1$ is admissible because every secondary
term above is $o(1)$ relative to the mains even at $X\asymp T$;
a reader preferring a power saving throughout may keep
$\lambda<1$ fixed and let $\lambda\to1^-$, with the same
limits \cite[Rem.~6.1]{C26}. The taper tail (the $2w/L$ deficit
in \eqref{e:windowcalc}) vanishes as $w/L\to0$; all endpoint
statements in this paper take the double limit in this order
(cf.\ \cite[\S7.1]{C26}).
\end{remark}

By Proposition~\ref{p:mu12} and the ledger of \S\ref{s:cert},
in the endpoint regime,
\begin{equation}\label{e:moments}
\mu_1\to1,\qquad \mu_2\to\tfrac43,\qquad \mu_3\to2,\qquad
\mu_4=\tfrac{197}{60}+R+o(1),
\end{equation}
where $R$ is the shifted-correlation class (\S\ref{s:cert:R}),
of model value exactly $-\tfrac1{30}=\tfrac{13}4-\tfrac{197}{60}$.
The zero-side consumption is the Christoffel-function count of
Proposition~\ref{p:count}: a lower bound on the number of
eigenvalues of $\widetilde G$ exceeding the vanishing threshold
$\theta_0$ of Proposition~\ref{p:tail} converts into the zero
counts \eqref{e:count}, and the mass bound near the origin comes
from the moment data alone (\S\ref{s:cert}). At the moment sequence
$(1,1,\tfrac43,2,\tfrac{13}4)$ the count gives $13/18$ and
$31/36$ exactly; it is monotone in each even moment and
decreasing in the odd-moment deficits.

Windows and cells: windows $I_m=[X,4X]$ and $I_n=(b_2/b_1)I_m$ of
length $Y\asymp X$ (more generally $Y\asymp X^{\vartheta}$,
$\vartheta\in(\tfrac12,1]$); prime-power moduli $b_1,b_2$ with
$g=(b_1,b_2)$, $r=b_1/g$, $q=b_2/g$; shifts $k\le K\asymp
Y/\max(r,q)$ and locks $|j|\le J\asymp X$; $\e(x)=e^{2\pi ix}$;
windowed prime sums on the position lattices
\[
S_m(\beta)=\sum_{m\in I_m}\LL(m)\,\e(b_2m\beta),
\qquad
S_n(\beta)=\sum_{n\in I_n}\LL(n)\,\e(b_1n\beta).
\]
$N_4(k,j)$ denotes the $\LL^4$-weighted count of locked
quadruples $(m,\,m-rk,\,n,\,n-qk)$ with $b_1n-b_2m=j$, and
$\SSS_4(k,j)\,V(k,j)$ its singular-series model, where
$V(k,j)$ is the window-overlap volume of the cell --- the
archimedean count of quadruples in the joint window under the
lock constraint, of size $\asymp Y^{2}/(rq\,\ell_1)$ after
consumption normalization --- with the conventions of the
fourth-moment cell ledger (\S\ref{s:fm:R}), frozen in the
reproduction suite. The five-case local table
$\kappa_p(v_p(j);v_p(b_1),v_p(b_2))$ defining
$\SSS_{b_1,b_2}(j)=\prod_p\kappa_p$ is verified digit-exact
against residue counting, with $\mathbb E_j[\SSS]=1$
(\S\ref{s:fm:toolkit}).

\emph{Import status.} \cite{C26} is a publicly posted
preprint; its framework is consumed exactly as stated in this
section, and its own analytic imports are the published
literature \cite{Wei52,IK,Tit86,Mon73,BGST}. The averaged
shifted-convolution input of the one-sided route is
Matom\"aki--Radziwi\l{}\l--Tao \cite{MRT}, used as
published. The formalization status of the paper's layers is
graded in \S\ref{s:verif}(f): the identity, local-law and
certificate layers are kernel-compiled in core Lean~4; the
analytic layer is not formalized.

\section{The certificate layer}\label{s:cert}

This section presents the certificate that carries the ladder
beyond the rank--trace inequality of \cite[Lem.~3.2]{C26}, and
the program's unconditional moment bookkeeping, machine-verified
throughout. It supplies: the Chebyshev--Markov certificate with
its exact $\kappa(c)$ laws; the third-moment theorem; the
$197/60$ decomposition of the fourth moment with the
two-bookkeeping reconciliation; and the two-wall structure of the
remaining class $R$, with its audit apparatus.

\subsection{The Chebyshev--Markov certificate and the exact
exchange rates}

\begin{lemma}[Certificate; machine-verified exactly]\label{l:cert}
Let $\nu$ be a probability measure on $\mathbb R$ with moments
$m_k$. Put
\[
Q(x)\;=\;\frac{\bigl(x^2-\tfrac{21}{8}x+\tfrac32\bigr)^2}{(3/2)^2}.
\]
Then $Q\ge0$, $Q(0)=1$, $Q\ge1$ on $(-\infty,0]$, $Q$ is
decreasing on $[0,\tfrac{21}{16}]$, and
\begin{equation}\label{e:EQ}
\int Q\,d\nu \;=\; \frac{4}{9}\Bigl(m_4-\tfrac{47}{16}\Bigr)
\;+\;R_Q,\qquad
R_Q:=\tfrac{4}{9}\Bigl[-\tfrac{21}{4}(m_3-2)
+\tfrac{633}{64}\bigl(m_2-\tfrac43\bigr)
-\tfrac{63}{8}\,(m_1-1)\Bigr],
\end{equation}
so $|R_Q|\le11\varepsilon$ if $|m_1-1|,|m_2-\tfrac43|,
|m_3-2|\le\varepsilon\le1$. For $0<\theta\le\tfrac1{100}$,
\[
\nu\bigl((-\infty,\theta]\bigr)
\;\le\;\bigl(1+4\theta\bigr)
\Bigl[\tfrac{4}{9}\bigl(m_4-\tfrac{47}{16}\bigr)
+11\varepsilon\Bigr].
\]
If $m_1=1$, $m_2=\tfrac43$, $m_3=2$, $m_4\le\tfrac{13}4$, then
$\nu((-\infty,\theta])\le\tfrac5{36}+O(\theta)$, and this is
sharp: the three-atom measure supported on $\{0,a,b\}$, with
$a=0.840635\ldots$, $b=1.784365\ldots$ the roots of
$x^2-\tfrac{21}8x+\tfrac32$, attains it. Consequently, under
$m_4\le\tfrac{13}4+c$ the counts of \S\ref{s:frame} give the
exact linear laws
\[
\liminf\frac{N_0^{s}}{N}\;\ge\;\frac{13}{18}-\frac{8c}{9},
\qquad
\liminf\frac{N_d}{N}\;\ge\;\frac{31}{36}-\frac{4c}{9}.
\]
\end{lemma}

\begin{proof}
Nonnegativity is clear and $Q(0)=1$; for $x\le0$ the inner
quadratic is $\ge\tfrac32$, so $Q\ge1$; on $[0,a]$ the inner
quadratic decreases from $\tfrac32$ through $0$ and $a>\tfrac1{100}$;
a direct computation gives $Q(\theta)^{-1}\le1+4\theta$ for
$\theta\le\tfrac1{100}$. Expanding $Q$ and taking expectations,
\[
\Bigl(\tfrac32\Bigr)^2\!\int Q\,d\nu
=m_4-\tfrac{21}4\,m_3+\Bigl(\tfrac{441}{64}+3\Bigr)m_2
-\tfrac{63}8\,m_1+\tfrac94,
\]
and substituting $m=(1,\tfrac43,2,m_4)$ gives
$\tfrac49(m_4-\tfrac{47}{16})$; the displayed $R_Q$ collects the
first-order deviations, and
$\tfrac49(\tfrac{21}4+\tfrac{633}{64}+\tfrac{63}8)\le11$. For
sharpness, the three-atom measure with moments
$(1,\tfrac43,2,\tfrac{13}4)$ exists by the Hankel recurrence
$m_{k+1}=\tfrac{21}8m_k-\tfrac32m_{k-1}$ on the atom set, and $Q$
vanishes doubly at $a,b$. The linear laws follow by inserting the
mass bound into the counts of \S\ref{s:frame} exactly as in the
proof display of \eqref{e:EQ}: mass
$\le\tfrac5{36}+\tfrac{4c}9+O(\theta)$, and
$1-2(\tfrac5{36}+\tfrac{4c}9)=\tfrac{13}{18}-\tfrac{8c}9$,
$1-(\tfrac5{36}+\tfrac{4c}9)=\tfrac{31}{36}-\tfrac{4c}9$.
\end{proof}

\begin{remark}\label{r:certlp}
Every identity of Lemma~\ref{l:cert} is verified symbolically
(sympy, exact rationals) in the reproduction package
(\texttt{certificate\_verification.py}) and re-run in the
certification audit; the roots, the $5/36$ evaluation, the $\kappa(c)$
laws, and the Hankel floors $m_5^{\mathrm{ext}}=\tfrac{177}{32}
=5.53125$, $m_6^{\mathrm{ext}}=\tfrac{2469}{256}=9.6445\ldots$
reproduce exactly. Re-optimizing the certificate polynomial for
each $c$ (a linear program) improves the record-breaking
threshold from $c<0.0559$ to $c<0.0721$
(\texttt{moment\_lp\_reopt.py}: $M^{*}=3.3221=\tfrac{13}4
\times1.0222$); the moment-pinned LP of \S\ref{s:lp} is the same
mechanism run at the full moment sequence.
\end{remark}

\subsection{The third moment, unconditionally}\label{s:cert:m3}

\begin{lemma}[Exact $3$-cycle resonances]\label{l:m3res}
Let $\lambda\le1$ and let $n_1,n_2,n_3\le X=T^{\lambda}$ be prime
powers with $|\log(n_1n_2/n_3)|\le C/T$. Then for $T\ge T_0(C)$,
$n_1n_2=n_3$; the analogous statement holds for every sign
pattern of the $3$-cycle.
\end{lemma}

\begin{proof}
$n_1n_2$ and $n_3$ are positive integers with ratio $e^{O(1/T)}$;
if distinct, their ratio differs from $1$ by at least
$1/n_3\ge T^{-\lambda}$, incompatible for large $T$ unless
$\lambda=1$; the borderline is handled by $|n_1n_2-n_3|\ge1$
versus $\le Cn_3/T\le C$: for $n_3\le T/(2C)$ equality is again
forced, and the range $n_3\in(T/2C,T]$ contributes only through
the taper tail, which vanishes in the normalized limit
(Remark~\ref{r:endpoint}).
\end{proof}

\begin{theorem}[Third moment]\label{t:m3}
For every fixed $\lambda\in(0,1]$ and admissible taper, $\mu_3$
converges as $T\to\infty$, and its limit tends to $2$ in the
endpoint limit $\lambda\to1^-$, $\eta\to0^+$.
\end{theorem}

\begin{proof}[Proof (nine-class ledger; full details in the
archived memo \texttt{V1\_completion.md})]
Expand $\operatorname{tr}\widetilde G^{3}$ against the $3$-cycle
$\Phi(\tau_1-\tau_2)\Phi(\tau_2-\tau_3)\Phi(\tau_3-\tau_1)$ after
the Gabor--Poisson collapse (Lemma~\ref{l:gabor}); the tail
correction is controlled by the trace-H\"older splitting
$\operatorname{tr}(A+E)^3-\operatorname{tr}A^3$ with
$\|E\|_{\mathrm{op}}\le\theta_0\ll lT^{\lambda/2-1}$
(Proposition~\ref{p:tail}) and $\|A\|_F^2=O(d\ell_1^2)$:
$o(d\ell_1^3)$ in total. Of the nine classes of
$\nu_X^{\otimes3}$, $\nu_X=\mu+\Pi_X+P_X$: (i) every class
containing $\Pi_X$ is $O(T^{\lambda/2-1})$ pointwise on the
window, negligible; (ii) single-$P$ classes are
$O(\sqrt X\,l^{2}L)$ by $\sum\LL(n)n^{-1/2}\le3\sqrt X$; (iii)
the class $\mu^3$ equals $d\ell_1^3(1+O(1/l))$ (the
$\ell_1$-normalization is exact to $O(1/l)$; measured
$1.0048/1.0014/1.0004$ at $T=600/38400/10^{8}$); (iv) each of
the three $\mu P^2$ classes reduces, by the $k=2$ case of
Lemma~\ref{l:m3res} plus the Montgomery--Vaughan Hilbert
inequality for the kernel tails, to the exact diagonal
$\sum_{n\le X}(\LL(n)^2/n)\,w(\log n)$ with the same
$(L-u)_+$-shaped edge weight as at $k=2$ (structure constants of
the cycle expansion are local to edges and slots, hence pinned by
the $k=2$ case, where the total is Montgomery's evaluation), and
$\sum_{n\le X}(\LL^2/n)(L-\log n)=L^3/6\,(1+O(1/L))$: each class
contributes $\tfrac13d\ell_1^3(1+o(1))$ at the endpoint, the
three positions summing to $1$; (v) the class $P^3$ is forced by
Lemma~\ref{l:m3res} onto exact relations $n_1n_2=n_3$, i.e.\
powers of a single prime, and $\sum_p\sum_{a+b=c}\log^3p\,p^{-c}
=O(1)$ makes it $O(Tl)$. Summing: $\mu_3\to1+1=2$. Cross-checks:
the $k=2$ instance of the same constant chain reproduces
Montgomery's $\tfrac13$ to four digits; the model gate
$\int O_3C_3=0$ is exact (\S\ref{s:cert:model}); the
finite-height shape $m_3m_1/m_2^2$ matches the model to $2\%$.
Two nonstandard steps --- the $k=3$ tail adaptation of
Proposition~\ref{p:tail} and the locality-plus-calibration fixing of
the constant chain --- are written out in the archived memo and
are consumed at that rigor.
\end{proof}

\subsection{The fourth moment: the $197/60$ ledger and the class
$R$}\label{s:cert:R}

\begin{proposition}[Evaluable classes of the fourth
moment]\label{p:m4classical}
By the identical ledger at $k=4$ ($81$ classes; $\Pi_X$- and
single-$P$ classes negligible; paired sign patterns forced exact
by Lemma~\ref{l:m3res}; $\{2,2\}$-pairing weights explicit window
integrals),
\[
\mu_4\;=\;\underbrace{1}_{\mu^4}
\;+\;\underbrace{2}_{\mu^2P^2\ \mathrm{diag}}
\;+\;\underbrace{\tfrac{17}{60}}_{\{2,2\}\ \mathrm{bare\
pairings}}\;+\;R\;+\;o(1)\;=\;\tfrac{197}{60}+R+o(1)
\]
in the endpoint regime, where $\tfrac{17}{60}=2\kappa$ with
$\kappa=\int_{[0,1]^2}xy\,\Sigma_{\mathrm{pair}}O_4
=\tfrac{17}{120}$ (exact quadrature of the certified $4$-cycle
overlap geometry), and $R$ is the shifted-correlation class:
pairs of restricted products $N_1=b_1m_1\ne N_2=b_2m_2$ (all legs
prime powers $\le X$), weight $\LL^4/\sqrt{N_1N_2}$, window
kernel $\widehat W(\Delta)=[\sin2T\Delta-\sin T\Delta]/\Delta$ at
$\Delta=\log N_1/N_2$, and the explicit overlap geometry $O_4$
(row/column asymmetric; unordered-cell weight
$Q(\beta_1,\beta_2)+Q(\beta_2,\beta_1)$); $R$-units
$=(1/\pi)\Sigma/(d\ell_1^4)$. The model value of $R$ is exactly
$-\tfrac1{30}=\tfrac{13}4-\tfrac{197}{60}$.
\end{proposition}

\begin{remark}[Bookkeeping reconciliation; audit
rounds 25--26]\label{r:books}
Two bookkeeping conventions coexist in the program archive and
must not be mixed. The \emph{bare-diagonal} convention above
books the $\{2,2\}$ pairings at $\tfrac{17}{60}$ and carries the
pair-level Montgomery plateau (the $\min(u,1)$ saturation of
$C_2$) inside $R$. The \emph{sine-model cumulant} convention
books $\{2,2\}$ as $2t_{\mathrm{adj}}+t_{\mathrm{opp}}$ with the
connected part separate. With the anchors now certified as exact
rationals (\S\ref{s:five:cert}: $t_{\mathrm{adj}}=\tfrac7{60}$,
$t_{\mathrm{opp}}=\tfrac1{30}$), the reconciliation is exact and
symmetric:
\[
\underbrace{\tfrac{17}{60}}_{\text{bare}}
=\underbrace{\tfrac4{15}}_{2t_{\mathrm{adj}}+t_{\mathrm{opp}}}
+\underbrace{\tfrac1{60}}_{\text{plateau}},
\qquad
R_{\mathrm{model}}
=\underbrace{-\tfrac1{60}}_{\text{plateau (pair-of-pairs Wick)}}
+\underbrace{-\tfrac1{60}}_{\text{connected core}}
=-\tfrac1{30},
\]
and $1+2+\tfrac{17}{60}=\tfrac{197}{60}$,
$2\cdot\tfrac{17}{120}=\tfrac{17}{60}$ exactly. The archived
grid values ($0.2677$, plateau $0.0156$, connected $-0.0177$;
cross-convention measurement $-0.01567$ in
\texttt{m1\_suite.py}) carried a small grid bias in
$t_{\mathrm{adj}}$ (quoted $0.11717$, exact $\tfrac7{60}
=0.11667$; register item D7); the exact connected model value
$-\tfrac1{60}=-0.01667$ agrees with the cumulant engine's
independent grid measurement $-0.01663$ to $0.2\%$. The totals
consumed anywhere in this paper were and remain exact
($\tfrac{197}{60}$, $-\tfrac1{30}$, $\tfrac{13}4$); the split is
now exact too. All ledgers below state which convention they use.
\end{remark}

\subsection{The two walls of $R$, and its direct
measurement}\label{s:cert:walls}

Everything in this subsection is unconditional structure or
finite-height measurement; it is the audit apparatus against
which the proofs of \S\S\ref{s:trunc}--\ref{s:four} were
developed and checked. (i)~\emph{Two ghost lemmas.} The
mean-field quadruple term is $O(l^{-3})$: the $u$-Jacobian
cancels $1/\sqrt{N_1N_2}$ exactly and the inner kernel integral
is $-2[\Si(2T/N)-\Si(T/N)]$, a boundary layer at $N\asymp T$; and
the smooth part of the cluster-squared sector vanishes
identically, since $\int_{\mathbb R}\widehat W_{\mathrm{in}}\,
d\Delta=0$. (ii)~\emph{Two walls.} Consequently
$R=R_{\mathrm{pp}}+R_{\mathrm{conn}}$: an edge layer carried by
Poisson-dual lattice terms at top scale (some leg within $O(1)$
of $X$; the explicit-formula regime) and a bulk object (model
bands at $s\sim1.5$). (iii)~\emph{Separation observable,
certified at three heights}: binning by $e=l-\max u_i$ gives, at
$l=5.95/6.64/7.33$: edge $-0.0061/-0.0068/-0.0060$ (saturating
toward the model $-0.0156$) and bulk $+0.0004/-0.0012/-0.0015$
(sign turn at $l\approx6$, monotone toward $-0.0177$).
(iv)~\emph{Direct measurement.} Meet-in-middle enumeration of all
resonant quartic prime pairs (exact kernels, class-resolved, flat
window, $\lambda=1$) gives
\[
\begin{array}{lcccc}
l & 4.56 & 5.94 & 7.33 & 8.72\\
R\ (\text{per }d\ell_1^4) & -0.0008 & -0.0057 & -0.0074 & -0.0101
\end{array}
\]
--- negative at every height, trending at the $1/l$-rate of all
calibration gates; a two-code-path audit showed the sharp cutoff
$C=40$ understates $|R|$ (full-kernel $-0.0012$ and $-0.0098$ at
the first two heights), so these are magnitude lower bounds. The
same audit certified the ledger's constant chain end-to-end at
$k=2,3,4$ to $0.05$--$0.8\%$. The $k{=}1$ Poisson stationary term
of the edge layer is evaluated in closed form (zone-1 master
formula; validated at three heights and two tapers to
$5$--$18\%$): a positive decaying transient $\sim c\,l/\ell_1^4$
explaining the late sign turn of the near band.

\subsection{Model calibration}\label{s:cert:model}

The free-fermion cumulant engine computes the Fourier-field
cumulants of the sine process by the ordered-set-partition
expansion of $\operatorname{Tr}\log(1+(e^f-1)K)$; gates:
$C_2(v)=\min(|v|,1)$ exactly, $\int O_3C_3=0$ (whence $m_3(1)=2$),
and the $m_4(1)$ assembly $3.25103$ against $\tfrac{13}4$ (grid
error $0.03\%$). The value $m_4(1)=\tfrac{13}4$ is pinned three
independent ways (cumulant assembly; GUE simulation
$m_{1..4}=1.0000,1.3313,1.9913,3.2213(27)$; Hankel forcing from
the extremal constants). GUE double-window Richardson
extrapolation measures the model values
\begin{equation}\label{e:guebands}
m_5(1)=5.573\pm0.113,\qquad m_6(1)=10.02\pm0.27,
\end{equation}
(sine-model/GUE measurement --- see the discrepancy register,
Appendix~\ref{a:record}, item D2, for the provenance correction
of an earlier attribution), and the six-moment certificate at the
central values gives $0.7328/0.8664$: the ladder climbs only
through the odd moments' strict-positivity gaps, since the
$k\le4$ extremal already has $m_5=\tfrac{177}{32}$,
$m_6=9.6445$. The assembled values of \S\ref{s:five} lie inside
the bands \eqref{e:guebands}.

\section{The one-sided fourth-moment route:
$0.6916$}\label{s:fm}

This section establishes a one-sided bound on $R$ sufficient to
cross the pair-correlation ceiling $0.68185$, by classical
technology plus the averaged shifted-convolution theorem of
\cite{MRT}. It is logically
independent of the exact evaluation of
\S\S\ref{s:trunc}--\ref{s:four} and is retained both as the
program's first crossing of the ceiling and because its
architecture (cells, bridge, aggregation, discharge ledger) is
consumed by the Lemma D proof.

\begin{theorem}[One-sided fourth-moment bound; grade:
certified-candidate with the numerical-analysis step of
Lemma~\ref{l:CL}]\label{t:fm}
With $R=R(1)$ the class of Proposition~\ref{p:m4classical},
$\limsup_{T\to\infty}R(1)\le\bar R:=0.0066$. Consequently, by the
count of \S\ref{s:frame} at
$(1,1,\tfrac43,2,\tfrac{13}4+\Delta)$, $\Delta=\bar R+\tfrac1{30}
=0.039933$,
\[
\liminf\frac{N_0^{s}}{N}\ \ge\ 1-2\Lambda_2(0)=0.691615,
\qquad
\liminf\frac{N_d}{N}\ \ge\ 1-\Lambda_2(0)=0.845807,
\]
where $\Lambda_2(0)=\dfrac{5/108+\Delta/3}{1/3+4\Delta/3}
=0.154193$, and the same for any fixed primitive Dirichlet
$L$-function. Any $\bar R\le0.0222$ crosses the $0.68185$
ceiling; the margin is $0.0156$ in $R$-units.
\end{theorem}

\noindent(An earlier draft printed $\Lambda_2(0)=0.156293$;
the certification audit recomputed the closed form and found the
stated theorem constants $0.6916/0.8458$ correct and that figure
stale --- discrepancy register D1, Appendix~\ref{a:record}.) The
constant $\bar R$ decomposes as $[\mathrm{core}\le-0.0055]+
[\varepsilon_{\mathrm{tail}}\le0.0111]+[\varepsilon_{CL}=0.0010]$,
where $\varepsilon_{CL}$ is a convergence-slack allowance for the
core sequence (Lemma~\ref{l:CL}).

\subsection{Toolkit: singular series, arc identity,
tails}\label{s:fm:toolkit}

For $b_1,b_2\ge1$, $j\ne0$ the local density of
$b_1m_1-b_2m_2=j$ is $\SSS_{b_1,b_2}(j)=\prod_p\kappa_p$ with the
five-case table verified digit-exact; $\mathbb E_j[\SSS]=1$. Its
arc expansion:

\begin{proposition}[Arc identity]\label{p:S}
$\displaystyle
\SSS_{b_1,b_2}(j)=\sum_{q\ge1}\beta_q\,c_q(j),\qquad
\beta_q=\frac{\mu(q_1)\mu(q_2)}{\varphi(q_1)\varphi(q_2)},\quad
q_i=\frac{q}{(q,b_i)}.$
\end{proposition}

\noindent At $b_1=b_2=1$ this is the classical evaluation used in
\cite[(66)--(67)]{MRT}; the general case follows by the same
$p$-local comparison against the $\kappa_p$ table. The truncation
tail obeys, for $g=(b_1,b_2)\le(\log M)^{C}$,
\begin{equation}\label{e:X2}
\textstyle\sum_{|j|\le H}\bigl(\SSS-\SSS^{(P)}\bigr)^2(j)\;\ll\;
H\,(\log)^{-B+O(1)},\qquad P=(\log)^{B},
\end{equation}
by the $(q_1,q_2)$-grading and $\sum_{|j|\le H}c_q(j)^2\ll(H+q)
\varphi(q)$. \emph{Discharge ledger} (kept, per the program's
audit discipline): cells with $g>(\log)^{C}$ are discarded at
certified cost $0.71\%$ weight share, $\le3\times10^{-4}$
$R$-units. The trivial envelope
$\sum_{n\le x}\LL(n)\LL(n+h)\ll\SSS(h)x$ uniformly for
$|h|\le x$ is \cite[Cor.~3.14]{HR}.

\subsection{Covered zone, middle band, bridge and
aggregation}\label{s:fm:R}

Write $S_i(\alpha)=\sum_{m\sim M_i}\LL(m)\e(b_im\alpha)$. The
two-modulus structure is discharged by four archive lemmas
(rounds 38--40): the major-arc law
$S_i(\alpha)=\frac{\mu(q_i)}{\varphi(q_i)}v_i(\eta)
+O(M_i\exp(-c\sqrt{\log M}))$ with $q_i=q/(q,b_i)$
($b$-uniformity free); the classification of joint near-major
points ($q\le\min(P^2g,Pb_1,Pb_2)$, a finite $b$-independent
family whose content is the $(q_1,q_2\le P)$-graded truncation
$\SSS^{(P)}$); the dilation transfer ($\gamma_i=b_i\alpha$ obeys
Vaughan's estimate in \emph{its own} approximation quality); and
a smoothing-collar lemma aligning the Gallagher width. After this
discharge the $b=1$ machinery of \cite{MRT} applies verbatim:

\begin{proposition}[{\cite[Thm.~1.3(i)]{MRT}}, variance
form]\label{p:IM}
For $X^{8/33+\varepsilon}\le H\le X^{1-\varepsilon}$, any center
$h_0\le X^{1-\varepsilon}$, and any $A$:
\[
\sum_{|h-h_0|\le H}\Bigl|\sum_{X<n\le2X}\LL(n)\LL(n+h)-\SSS(h)X
\Bigr|^2\;\ll_{A,\varepsilon}\;HX^2(\log X)^{-A}.
\]
\end{proposition}

\noindent The covered share of the class at the $8/33$ exponent
is $0.885$ (anchored at the archived $0.797$ at exponent
$\tfrac13$). For the middle band, with
$A(b_2,j)=\sum_m\LL(m)\LL((b_2m+j)/b_1)\mathbf1[b_1\mid b_2m+j]$
and $MT=\SSS_{b_1,b_2}(j)\,\mathrm{len}/(b_1b_2)$:

\begin{lemma}[Dispersion swaps; exact]\label{l:swaps}
$\sum_{b_2,j}A^2$, $\sum A\cdot MT$ and $\sum MT^2$ are, exactly,
weighted twin-sum families over the structured shifts $h=qk$
(outer pairs $m-m'=rk$), the $\SSS$-against-$\LL$ correlation on
the strip, and $2J\,E_2(b_1,b_2)(\mathrm{len}/b_1b_2)^2(1+o(1))$
respectively, where $E_2=\prod_p\mathbb E_v[\kappa_p^2]$ (anchor:
$E_2(1,1)=2\prod_{p>2}(1+(p-1)^{-3})$).
\end{lemma}

\begin{lemma}[Bridge]\label{l:bridge}
For any window $W$ and modulus $q$:
$\displaystyle\sum_{k\ne0}\ \sum_{n,n-qk\in W}\LL(n)\LL(n-qk)
=\sum_{a\bmod q}\Bigl[\psi_W(q,a)^2-\sum_{n\in W,\,n\equiv a}
\LL(n)^2\Bigr].$
\end{lemma}

\begin{lemma}[Structured-to-full aggregation]\label{l:U}
Let $W$ have length $Y$, $\mathrm{dev}(h)$ the (smoothly
weighted) twin deviation on $W$, $Q$ a set of moduli with shift
counts $K_q\le H/q$, $H\le Y^{1-\varepsilon}$. If
$\sum_{h\le H}|\mathrm{dev}|^2\ll_{A'}HY^2(\log)^{-A'}$ for every
$A'$, then for every $A$
\[
\sum_{q\in Q}\sum_{k\le K_q}|\mathrm{dev}(qk)|^2
\;\ll_A\;HY^2(\log Y)^{-A},
\]
uniformly in $Q$ --- no factor of the moduli is lost.
\end{lemma}

\begin{proof}
The left side is $\sum_{h\le H}\nu(h)|\mathrm{dev}(h)|^2$ with
$\nu(h)\le\tau(h)$. Split at $\nu\le(\log Y)^{C}$: the generic
part costs $(\log)^{C}$ against the arbitrary-$A'$ input; on the
rare part use the envelope $|\mathrm{dev}(h)|\le9\SSS(h)Y$ with
$\sum_{h\le H,\tau(h)>V}\tau(h)\ll H(\log)^3/V$ at
$V=(\log)^{C}$. Choose $C=A+4$, $A'=2A+4$.
\end{proof}

\noindent The band consumer applies a single Cauchy--Schwarz
\emph{in the shift variable only} (any Cauchy--Schwarz across the
family pays its Poisson floor, measured $50$--$60\times$ over
budget, and is never used). The glue layer (welding weight
$w_k(n)=\sum_{m\in I(n)}\LL(m)\LL(m-rk)$) closes by the exact
factorization of main terms, Abel summation against the glue for
the major arcs, and the divisor-bounded envelopes of
\cite[Prop.~5.4, App.~A]{MRT} for the minor arcs.

\subsection{Deterministic mains and the continuum
limit}\label{s:fm:mains}

Write $\SSS_{b_1,b_2}=\bar D+g_P+g_{\mathrm{tail}}$ ($\bar D$ the
exact finite local mean; $g_P$ the $q\le P$ oscillating piece;
$P=40$).

\begin{lemma}[Partial-sum identity]\label{l:PS}
For every $M\ge1$,
$\displaystyle\sum_{j\le M}(\SSS_{b_1,b_2}-\bar D)(j)=
-\sum_{d\ge1}d\,\bigl\{\tfrac Md\bigr\}\,
\gamma^{\mathrm{free}}_d$, with explicit Euler-local $\gamma$,
absolutely convergent against the $\ell^1$ weights
$d\min(1,M/d)$.
\end{lemma}

\begin{lemma}[Finite-$J$ evaluation and uniform
tail]\label{l:E3}
For every finite $J$ and $\theta$,
$\sum_{j\le J}\mathrm{tail}_P(j)\,\mathrm{row}_j(\theta)
=\sum_{d\ge1}\gamma_d^{>P}\,S_0(d\theta;\lfloor J/d\rfloor)$
with $S_0(x;K)=\sum_{k\le K}\frac{\sin2xk-\sin xk}{k}$; hence the
$\theta$-uniform bound $\le C_{\mathrm{saw}}\sum_d
|\gamma_d^{>P}|$. No $J\to\infty$ interchange is performed
anywhere.
\end{lemma}

\begin{lemma}[Certified sawtooth constant]\label{l:S0}
$|S_0(x;K)|\le C_{\mathrm{saw}}:=2.10$ for all $x$ and $K\ge1$.
\end{lemma}

\begin{proof}
For $K\le64$: rigorous grid verification with the Lipschitz bound
$|\partial_xS_0|\le3K$ certifies $\max|S_0|\le1.8675$. For
$K>64$: the Dirichlet representation with the sine-integral
comparison near the singular set and the second mean value
theorem away from it gives $\le2.05$; samples at $K=128,512,4096$
read $1.842,1.849,1.851$, approaching $\Si(\pi)=1.85194$ from
below. (Re-run in the certification audit: identical values.)
\end{proof}

\noindent Zone-integrating the uniform tail against the certified
geometry gives $\varepsilon_{\mathrm{tail}}\le0.0115$
($T=2400$), $0.0111$ ($T=4800$), decreasing; the evaluated core
(exact finite forms) is $\mathrm{core}(2400)=-0.00549$,
$\mathrm{core}(4800)=-0.0064$, certified against the pipeline
anchor ($-0.00545$ vs the suite's $-0.00544$).

\begin{lemma}[Continuum limit of the core; grade: proof shape
with computed constants --- see the grading note
below]\label{l:CL}
$\mathrm{core}(T)$ converges as $T\to\infty$ to
$C_{\mathrm{core}}=\sum_{b_1\le b_2}\frac{\LL(b_1)\LL(b_2)}
{b_1b_2}F_\infty(b_1,b_2)$, an absolutely convergent series of
computable continuum sawtooth integrals, with
$|\mathrm{core}(T)-C_{\mathrm{core}}|\ll(\log T)^{-1}$; the
ensemble collapses onto the universal $(1,1)$-class (Mertens mass
concentrates on pairs of large primes), whose slope and constant
are computed, not fitted: partial-sum slopes
$0.99684\to0.99962\to0.99996$ over successive decades (classical
law, exactly $1$ --- now \emph{proved}: the Dirichlet series is
$\zeta(s+1)E(s)$ with $2C_2E(0)=1$, Proposition~\ref{p:c0}),
and the constant $c_0=-1.4228$ is now \emph{identified in closed
form}: $c_0=\gamma-2=-1.4227843\ldots$ (Euler--Mascheroni
$\gamma$; Proposition~\ref{p:c0}, verified to $10^{-5}$ at
$M=2\times10^{7}$); $\theta$-ladder $c^{*}=-0.5433$,
decade-stable, matching the archive's independent $-0.55(2)$
($c^{*}=\log\pi+\gamma-2-\overline{\mathrm{osc}}$ in the pin of
the archived F-QC5 gate, so the remaining computed piece of
$c^{*}$ is the mean oscillation
$\overline{\mathrm{osc}}\approx0.2652$, not identified here). The continuum quadrature with Richardson
extrapolation in $1/l$ ($l=40,60,100,160$) yields
$C_{\mathrm{core}}=-0.0209\pm0.0026$, independently confirming
the archive ladder $-0.020(2)$ --- two disjoint routes at $4\%$.
\end{lemma}

\noindent\emph{Grading note.} The convergence structure
(scaled-variable Riemann sum, universality collapse) is an
analytic argument; the values of the collapsed series' slope and
constant, and the quadrature of $C_{\mathrm{core}}$, are
computed constants with stated bands, not certified enclosures.
Theorem~\ref{t:fm} charges only the conservative
$\mathrm{core}\le-0.0055+\varepsilon_{CL}$,
$\varepsilon_{CL}=0.0010$, against the measured decreasing grid;
the certification write-out of the quadrature is a named open
item (\S\ref{s:verif}). Substituting the computed
$C_{\mathrm{core}}$ would give $\bar R=-0.0098\pm0.0026<0$,
hence $0.7031/0.8515$; that upgrade is \emph{not} consumed
anywhere in this paper.

\subsection{Assembly}\label{s:fm:assembly}

At $\lambda=1$, uniform full-kernel convention:
\[
R(1)\;\le\;\mathrm{core}(T)+\varepsilon_{\mathrm{tail}}+o(1)
\;\le\;(-0.0055+0.0010)+0.0111+o(1)\;=\;0.0066+o(1),
\]
the $o(1)$ collecting the covered-zone, band and low-zone
deviations of \S\ref{s:fm:R} (the low-zone cap's certified
finite-height price at the operating point is $0.0163$, inside
the LP budget $0.0387+0.0177$). Feeding $\bar R=0.0066$ into the
count gives Theorem~\ref{t:fm}.\qed

\section{The truncation discharge}\label{s:trunc}

\begin{lemma}[Cell-mass bound]\label{l:mass}
Uniformly in the lock and shift ranges, the $\LL^{2k}$-weighted,
consumption-normalized mass of the $k$-cycle cells at parameter
$\ell_1$ is $\ll Y^{2}(\log Y)^{c}\,\ell_1$.
\end{lemma}

\begin{proof}
For a fixed cell, each locked configuration is determined by
$(m,k)$ together with the residue data of the lock; the count of
$(m,k)$ pairs in the windows is $\ll YK\ll Y^{2}/\max(r,q)$, each
carrying weight $\ll(\log Y)^{2k}$, and the lock constraint
selects a $1/\ell_1$-proportion of the $n$-range up to
divisor-bounded multiplicity; the consumption normalization
(division by the ledger volumes $V\asymp Y^{2}/(rq\ell_1)$)
leaves at most one factor $\ell_1$. The numerical ledger face
measures the aggregate mass essentially flat in $\ell_1$ (fitted
exponent $-0.11$), well inside the bound.
\end{proof}

\begin{lemma}[Multiplicity]\label{l:mult}
The number of unordered prime-power tuples with cell parameter
$\ell_1=\ell$ is $\ll_{\varepsilon}\ell^{\varepsilon}$
(divisor-bounded).
\end{lemma}

\begin{proposition}[Truncation--consumption
compatibility]\label{p:consume}
Let $F$ denote the consumed $k$-cycle cell aggregate, $k\ge4$,
and $F^{(P)}$ its restriction to $\ell_1\le P$. Then
\[
\bigl|F-F^{(P)}\bigr|\;\ll_{\varepsilon}\;
Y^{2}(\log Y)^{c}\,P^{-(k-2)+\varepsilon}
\qquad\text{two-sidedly.}
\]
Hence for every $A_0$ there is $B$ (any fixed $B>(c+A_0)/(k-2)$
suffices) such that proving the variance bound for moduli
$\le P=(\log X)^{B}$ proves it in the consumed form, with total
truncation cost $\ll(\log X)^{-A_0}$: only moduli
$\le(\log X)^{B}$ are ever consumed.
\end{proposition}

\begin{proof}
Weight $\times$ mass $\times$ multiplicity
$=\ell^{-k}\cdot\ell\cdot\ell^{\varepsilon}$ summed over
$\ell>P$. The bound is on absolute values, hence two-sided. The
archived D1 layer independently certifies the squared-remainder
grading at $k=4$ ($\mathbb E_j[(F-F^{(P)})^2]\ll P^{-1-\delta}$;
tail locals $t_p^{2}\sim4/p^{2}$). Two truncations coexist in the
program and must not be confused: the sawtooth/frequency
truncation inside the $C_{\mathrm{full}}$-quadrature machinery of
\S\ref{s:fm:mains}, whose tail carries that limit, and the
present $\ell_1$-cell truncation, whose tail is trivially small;
the audit trail (archive rounds 91--93, item S1) verified the
disentanglement. An earlier draft stated the multiplicity as
$\ell_1$; the divisor bound of Lemma~\ref{l:mult} is what the
displayed exponent uses.
\end{proof}

\section{The local closure}\label{s:local}

The local (mod $p$) structure of every moment order is governed
by affine lock chains. Fix a cycle length $b\ge2$, a prime power
$P_0=p^{s}$, unit coefficients $a_1,\dots,a_b$, and free locks
$j_1,\dots,j_{b-1}\in\mathbb Z/P_0$; define points by
$n_{t+1}\equiv a_tn_t+j_t$, the last lock induced by cyclic
closure. Write $n_t=A_t(n_1+d_t)$ where $A_{t+1}=a_tA_t$,
$A_1=1$, $d_1=0$; then $d_{t+1}-d_t=j_t/A_{t+1}$ is a unit
multiple of $j_t$. Let $N(j)$ count the $n_1\in\mathbb Z/P_0$
with all $b$ points prime to $p$.

\begin{proposition}[Coincidence-counting closure]\label{p:local}
Pointwise in the lock vector $j$,
\[
N(j)\;=\;P_0\,\frac{p-D(j)}{p},
\qquad
D(j)=\#\{d_t\bmod p:\ t=1,\dots,b\},
\]
and for any class $\mathcal C$ of locks,
\[
\kappa_b(\mathcal C)
=\frac{1-\mathbb E[D\mid\mathcal C]/p}{(1-1/p)^{b}},
\qquad
\mathbb E_j[\kappa_b]=1\ \text{exactly.}
\]
Forbidden-class merges are pair events of weight $1/(p-1)$;
triple merges enter at $O(p^{-2})$.
\end{proposition}

\begin{proof}
Since $A_t$ is a unit, $n_t=A_t(n_1+d_t)$ is prime to $p$ iff
$n_1\not\equiv-d_t\pmod p$; the all-unit condition excludes
exactly the $D(j)$ distinct residues $\{-d_t\}$, each of density
$1/p$ in $\mathbb Z/P_0$, giving the count. The class formula is
its average against $(1-1/p)^{b}$. For the mean-one identity:
under uniform locks the increments are independent uniform, so
$(d_2,\dots,d_b)$ is uniform on $(\mathbb Z/p)^{b-1}$ and
$\mathbb E[D]=p(1-(1-1/p)^{b})$, whence
$1-\mathbb E[D]/p=(1-1/p)^{b}$. A merge across a gap of lock-sum
$\sigma$ requires $p\mid\sigma$: impossible for a single unit
lock; for a sum of two units it occurs for exactly one value of
the second, weight $1/(p-1)$.
\end{proof}

Machine verification (pipeline, \texttt{face\_local\_closure};
re-run in the certification audit with identical output): zero pointwise
violations over full enumerations and closed-form match to
$2\times10^{-16}$ for every lock class at $b=4,5,6$,
$p\in\{5,7,11,13\}$, $s\in\{1,2\}$, with the mean-one identity
exact. The $\kappa_4$ closed forms of the fourth-moment table are
instances ($\kappa_4(0,0)=1-(4-\tfrac2{p-1})/p$). The proposition
supplies the singular-series combs at every cycle length and the
$j$-averaged factorization identities used below. (A first draft
of the corresponding Lean seed omitted the normalising inverse in
$d_t=B_tA_t^{-1}$; the exact-enumeration proxy fired --- registry
r103, the $29$th conversion --- and the seed was corrected:
Appendix~\ref{a:record}.)

\section{Lemma D and the exact fourth moment}\label{s:four}

The fourth-moment evaluation $m_4\to13/4$ closes the estimate
isolated by the program (``Lemma D'': the averaged, lock-resolved
quadruple-correlation variance) in its consumed, truncated form.
In its \emph{unrestricted} form --- moduli up to a power of $X$
--- it remains open, one structural step beyond the published
diagonal cases of the Barban--Davenport--Halberstam/dispersion
lineage. Two observations remove the difficulty from the consumed
form: the truncation discharge of \S\ref{s:trunc}, and the
spectral frame below.

\subsection{The frame and the Parseval identity}

By the archived, machine-verified bridge and glue identities
(\S\ref{s:fm:R}), the $(k,j)$-aggregated variance for a modulus
pair is carried by the \emph{dilated density} $\rho(\tau)$: with
$A_m,A_n$ the autocorrelations of the two position-lattice
sequences, on the analysis grid of period $N$ (a power of two
exceeding four times the largest position, so no correlation
wraps),
\[
\rho(\tau)\;=\;\sum_{t}A_m(t)\,A_n(t-\tau)
\;=\;\mathrm{irfft}\bigl(|S_m|^{2}|S_n|^{2}\bigr)(\tau),
\]
and the lock-variance obeys the \emph{exact} discrete Parseval
identity
\begin{equation}\label{e:pars}
\sum_{\tau=0}^{N-1}\rho(\tau)^{2}
=\frac1N\Bigl[\mathrm{dens}_0^{2}
+2\!\!\sum_{0<\beta<N/2}\!\!\mathrm{dens}_\beta^{2}
+\mathrm{dens}_{N/2}^{2}\Bigr],
\qquad \mathrm{dens}=|S_m|^{2}|S_n|^{2},
\end{equation}
verified bitwise in the pipeline ($10^{-16}$ relative; re-run in
the certification audit). The same identity applied to
$\rho-\mathrm{Model}$ turns the variance into the spectral
$L^{2}$-distance of $\mathrm{dens}$ from the model comb. We
classify the frequencies $\beta$ by the arc quality of the pair
$(\gamma_m,\gamma_n)=(b_2\beta,b_1\beta)\bmod1$ at cutoff
$P^{*}=(\log X)^{B^{*}}$: both major (MM), mixed, or both minor.

\subsection{The three leaves}

\begin{lemma}[MM leaf]\label{l:MM}
Let $\beta$ be a doubly major bin: $\gamma_m=a/Q+\eta$,
$Q\le P^{*}$, $|\eta|\le$ arc width, and likewise $\gamma_n$.
Then, for every $A'$,
\[
S_m(\beta)=\frac{\mu(Q)}{\varphi(Q)}\,I_m(\eta)
+O_{A',B^{*}}\!\bigl(Y(\log Y)^{-A'}\bigr),
\qquad
I_m(\eta)=\sum_{m\in I_m}\e(m\eta),
\]
uniformly, and the same for $S_n$; the constants are ineffective.
Consequently, on the MM bins,
$\bigl|\mathrm{dens}-\mathrm{comb}\bigr|
\ll\mathrm{main}\cdot Y(\log Y)^{-A'}+Y^{2}(\log Y)^{-2A'}$
pointwise, where $\mathrm{comb}$ is the product of the squared
main terms.
\end{lemma}

\begin{proof}
$S_m(\beta)=\sum_m\LL(m)\e(m\gamma_m)$ is a windowed single prime
sum at $\gamma_m=a/Q+\eta$. Splitting into progressions mod $Q$
and applying the Siegel--Walfisz theorem in the form
\cite[\S22]{Dav} (or \cite[Cor.~5.29]{IK}) with partial summation
in the window against the phase $\e(m\eta)$ gives the display;
the constants depend only on $(A',B^{*})$ and are ineffective
through Siegel's theorem. The pointwise comb comparison follows
by expanding the product. The per-bin face
(\texttt{face\_sw}) measures the absolutely normalised error
across six families: RMS $3.6\times10^{-3}$ at $T=2400$ falling
to $3.6\times10^{-4}$ at $T=153600$.
\end{proof}

\begin{remark}[Ineffectivity enters here]\label{r:ineffective}
This is the single point where ineffective constants enter the
chain: Siegel's theorem gives no computable $T_0(A')$, so every
theorem constant downstream is ineffective, while every
finite-height face measured in \S\ref{s:faces} carries
effective constants. The two kinds of statement are graded
separately throughout (see the closing remark of
\S\ref{s:faces}); nothing else in the paper --- the identities,
the certificate, the local law, the window calculus --- consumes
an ineffective input.
\end{remark}

\begin{lemma}[Minor and mixed leaves]\label{l:minor}
For $\gamma$ minor at cutoff $P^{*}$ (denominator range
$((\log X)^{B^{*}},Y(\log X)^{-B^{*}})$ by Dirichlet
approximation), $|S(\gamma)|^{2}\ll Y^{2}(\log Y)^{-B^{*}/2+c_0}$
by Vaughan's estimate. Hence the doubly-minor spectral mass
carries a $(\log)^{-B^{*}+2c_0}$ saving squared, and each mixed
bin at least the one-sided saving; the mixed and minor
contributions to the spectral $L^{2}$-distance are
$\ll(\log X)^{-B^{*}/2+c_1}$ relative to the natural size. The
size bookkeeping (well-spacing of image families and the hybrid
large-sieve/minor-sup stratum budget) is classical and is
face-checked at $24$--$31\times$ tightness
(\texttt{face\_arcs}).
\end{lemma}

\begin{lemma}[The comb is the singular series]\label{l:comb}
The spectral main comb of Lemma~\ref{l:MM} equals the truncated
singular-series model: the surviving frequencies are exactly the
CRT progression $t\equiv0\ (b_2)$, $t\equiv\tau\ (b_1)$ in the
$\tau$-frame, and the local factors match the
coincidence-counting closed form of Proposition~\ref{p:local};
the truncation tail of the singular series beyond $P^{*}$ is
$\ll(\log X)^{-B^{*}+c}$ by the standard
$\sum_{q>P^{*}}\mu^{2}(q)(q,h)/\varphi(q)^{2}$ bound. The
identification is displayed explicitly, with a worked modulus
pair, in Appendix~\ref{a:comb}; the $\tau$-resolved model tracks
the measured profile at $0.2$--$4\%$ across sixteen families and
eight heights (\texttt{face\_dispersion}).
\end{lemma}

\subsection{The theorem}

\begin{theorem}[Lemma D, truncated form]\label{t:D}
Let $B$ be fixed, $b_1,b_2\le(\log X)^{B}$ prime powers, windows
of length $Y\asymp X^{\vartheta}$, $\vartheta\in(\tfrac12,1]$.
Then for every $A>0$,
\[
\sum_{k\le K}\sum_{|j|\le J}
\bigl|N_4(k,j)-\SSS_4(k,j)V(k,j)\bigr|^{2}
\;\ll_{A,B}\;KJV^{2}(\log X)^{-A}.
\]
Consequently, with Proposition~\ref{p:consume}, $m_4\to13/4$
two-sidedly.
\end{theorem}

\begin{proof}
Choose $B^{*}=2(A+c_1+B)$ and apply \eqref{e:pars} to
$\rho-\mathrm{Model}$. The spectral $L^{2}$-distance splits over
the arc classification. On MM bins, Lemma~\ref{l:MM} with
$A'=A+2B^{*}$ gives a pointwise relative saving $(\log)^{-A'}$,
hence $\sum_{\mathrm{MM}}|\mathrm{dens}-\mathrm{comb}|^{2}
\ll(\log)^{-2A'+c}\sum\mathrm{main}^{2}$; by Lemma~\ref{l:comb}
the comb agrees with the model to $(\log)^{-B^{*}+c}$ relative,
absorbed by the choice of $B^{*}$. The mixed and doubly-minor
bins contribute $\ll(\log)^{-B^{*}/2+c_1}$ relative by
Lemma~\ref{l:minor}. The $(k,j)$-aggregation constants of the
bridge and glue layer are archived and certified
(\S\ref{s:fm:R}); the normalization match to $KJV^{2}$ is the
consumer interface there. All bounds are on absolute values,
hence two-sided; all constants depend only on $(A,B)$ and are
ineffective. Choosing $B^{*}$ then $A'$ in that order involves no
circularity.
\end{proof}

\begin{remark}[The progression-frame route; falsifier
r95]\label{rem:prog}
The archived reduction of the variance to a locked two-modulus
progression covariance suggests a shorter argument: expand both
progression errors in additive characters and evaluate by
Siegel--Walfisz. Executed at the moduli $(r,q)$ separately this
produces \emph{incomplete} character sums (the lock couples
residues modulo $r$ to residues modulo $q$, and the $a$-sum over
one modulus has no orthogonality against the other); the
expansion must be lifted to the joint modulus
$M\asymp rq\le(\log X)^{2B}$, where it becomes a complete-sum
argument parallel to the one above. An earlier draft presented
that sketch without the lift; it was replaced by the
spectral proof, in which the $\tau$-transform runs over the full
period and the orthogonality (Lemma~\ref{l:comb}) is exact. The
slip and its correction are falsifier r95 on the forward record.
\end{remark}

\begin{corollary}[The $13/18$ rung]\label{c:1318}
With ineffective constants,
\[
\liminf_{T\to\infty}\frac{N_0^{s}(T,2T)}{N(T,2T)}\ge\frac{13}{18},
\qquad
\liminf_{T\to\infty}\frac{N_d(T,2T)}{N(T,2T)}\ge\frac{31}{36},
\]
and the same for any fixed primitive Dirichlet $L$-function.
\end{corollary}

\begin{proof}
Theorem~\ref{t:D} and Proposition~\ref{p:consume} give
$m_4=13/4+o(1)$ two-sidedly; Theorem~\ref{t:m3} gives $\mu_3\to2$;
the certificate (Lemma~\ref{l:cert}, equivalently the
Christoffel-function count of Proposition~\ref{p:count} at
$(1,1,\tfrac43,2,\tfrac{13}4)$) gives $13/18$ and $31/36$
exactly, monotonicity in the fourth moment converting
$m_4\to13/4$ into the two liminf bounds. The Dirichlet
$L$-function case is as in \cite[Thm.~E]{C26}: the framework
uses only the functional-equation symmetry
$\rho\mapsto1-\bar\rho$, the zero-density bound, and the
explicit formula, all available for any fixed primitive
Dirichlet $L$-function. Constants are
ineffective through Lemma~\ref{l:MM}.
\end{proof}

\section{The fifth and sixth moments}\label{s:five}

\subsection{The product identities}

\begin{proposition}\label{p:prod}
For a $b$-cycle of position lattices sharing a window, with any
per-lattice densities $f_i,g_i$:
$\sum_jN_f(j)N_g(j)=\sum_t\prod_iC_i(t)$,
$C_i(t)=\sum_xf_i(x)g_i(x+t)$. In particular the lock variance is
the pointwise product of the lattice autocorrelations, summed
over the common shift.
\end{proposition}

\begin{proof}
Two chains share a lock vector iff they differ by a common shift;
the chain sum then factorizes over the positions.
\end{proof}

This identity ($O(N\log N)$ to evaluate) transports the spectral
argument of \S\ref{s:four} \emph{factor by factor} to every cycle
length: each factor is a windowed single prime sum, so the
three-leaf closure (per-factor Siegel--Walfisz on all-major bins,
Vaughan on any minor factor, comb by Proposition~\ref{p:local},
truncation by Proposition~\ref{p:consume} with stronger
exponents) applies verbatim. Machine faces: the
twin-singular-series product model tracks the empirical five-fold
product at $0.28$--$2.5\%$ (six modulus patterns, $T=4800$ to
$76800$, improving with height) and the six-fold product at
$0.11$--$1.29\%$.

\subsection{Pair layers collapse onto certified anchors}

The Bell$(5)=52$ and Bell$(6)=203$ class ledgers are
machine-enumerated (class tables $1/10/15/10/10/5/1$ and
$1/15/45/15/20/60/10/15/15/6/1$; pentagon chords $10+5$; hexagon
chords $30+15$; matching crossing spectrum $5/6/3/1$).
Frozen-slot identities --- a singleton block repeats a walk value
and leaves the overlap range unchanged; machine-exact --- reduce
every pair class to the four-cycle geometry, whose anchors are
now \emph{exact rationals} (\S\ref{s:five:cert}):
$t_{\mathrm{adj}}=\tfrac7{60}$, $t_{\mathrm{opp}}=\tfrac1{30}$,
$\Phi_4=\tfrac14-2t_{\mathrm{adj}}-t_{\mathrm{opp}}=-\tfrac1{60}$
(sine-model cumulant convention, cf.\ Remark~\ref{r:books}).
Odd-block classes vanish by the local law. Hence
\[
m_5=\underbrace{1+\tfrac{10}3+10t_{\mathrm{adj}}
+5t_{\mathrm{opp}}+5\Phi_4}_{=\,67/12\ \text{exactly}}+\ C_5,
\qquad
m_6=\tfrac{39}4+6C_5+\{2,2,2\}+\{4,2\}+\{6\},
\]
the rational parts being anchor-free (the $t$-geometry cancels:
$10t_{\mathrm{adj}}+5\Phi_4+5t_{\mathrm{opp}}=\tfrac54$ and
$30t_{\mathrm{adj}}+15t_{\mathrm{opp}}+15\Phi_4=\tfrac{15}4$
identically --- verified symbolically in the certification
audit and kernel-compiled in the Lean package,
\S\ref{s:verif}(f)), with $\{2,2,2\}=\tfrac{131}{420}$ now
\emph{certified exactly} (Proposition~\ref{p:t222}),
$\{4,2\}=-0.0552(8)$ (the four-cycle connected object consumed
jointly with a spectator pair), and $\{6\}=-0.0078(8)$
(grouped midpoint-grid ladders, \S\ref{s:conv}). The former write-out items W1--W3
are written out in \S\ref{s:writeouts}.

\subsection{The pairing layer certified exactly}\label{s:five:cert}

\begin{proposition}[Exact pairing constants]\label{p:t222}
The four-cycle anchors and the three-chord classes are the exact
rationals
\[
t_{\mathrm{adj}}=\tfrac7{60},\quad
t_{\mathrm{opp}}=\tfrac1{30},\quad
T_0=\tfrac3{70},\quad T_1=\tfrac1{90},\quad
T_2=\tfrac1{180},\quad T_3=\tfrac1{70},
\]
whence $\{2,2,2\}=5T_0+6T_1+3T_2+T_3=\tfrac{131}{420}$
(Touchard--Riordan multiplicities $5/6/3/1$).
\end{proposition}

\begin{proof}
Each constant is the integral over a box of
$(1-\mathrm{spread})_+\cdot\prod\min(|x_i|,1)$ for an explicit
walk (the chord-pattern walks of the Bell ledgers): a
piecewise-polynomial integrand all of whose kink hypersurfaces
are $\{0,\pm1\}$-combinations of the variables at levels
$\{0,\pm1\}$. Since $0$ is a walk position and every variable is
a position or a difference of two positions, the support lies in
$[-1,1]^d$ and the $\min$-caps never bind. The integral is then
computed \emph{exactly} in rational arithmetic by iterated
integration with full breakpoint enumeration: inner breakpoints
have unit leading coefficients, so all collision and boundary
resultants remain small-integer forms, and on every kink-free
piece the integrand is a polynomial of degree at most $2/4/8$ (in
the inner/middle/outer variable), integrated exactly by closed
Newton--Cotes rules of order $5/7/9$ at rational nodes. Every
piece is integrated twice (whole piece, and the two halves): for
polynomials of degree below the rule order the two exact
rationals must coincide, and any missed kink forces a
discrepancy; no discrepancy occurred. $t_{\mathrm{opp}}$ and
$T_0$ additionally admit short closed-form evaluations by orthant
decomposition (Appendix~\ref{a:cert}), agreeing with the
algorithmic values. Code and outputs:
\texttt{certification/exact\_t222.py}.
\end{proof}

\begin{remark}[The anchor correction]\label{r:anchor}
The archived anchor $t_{\mathrm{adj}}=0.11717$ (rounds 69--71
grid) is corrected by Proposition~\ref{p:t222} to
$\tfrac7{60}=0.11\overline{6}$: a $+0.43\%$ rectangle-rule bias
at kinks. Nothing consumed downstream changes: the $m_5$ and
$m_6$ rational layers are anchor-free (above), and the
reconciliation of Remark~\ref{r:books} becomes exact. The
formerly conjectural rational identifications of the three-chord
classes are now theorems. Register items D7, D8.
\end{remark}

\subsection{The write-outs W1--W3}\label{s:writeouts}

The three items carried as named write-outs in the audit ledger
are discharged here; their pre-registered faces are in
\texttt{certification/writeouts\_verification.py} (all green;
one face fired during construction and was converted ---
Appendix~\ref{a:record}).

\begin{lemma}[W2: free-slot factorization]\label{l:W2}
In any $b$-cycle cell ledger, a slot whose lock is unconstrained
factorizes exactly: the cell aggregate equals the contracted
$(b-1)$-cycle aggregate times the window mass
$\sum_{m\in W}\LL(m)$ of the free slot, and
$\sum_{m\in W}\LL(m)=|W|\bigl(1+O(e^{-c\sqrt{\log X}})\bigr)$.
\end{lemma}

\begin{proof}
The lock sum over the free slot's offset is unrestricted, so the
double sum factorizes by Fubini; on the walk-geometry side this
is the frozen-slot identity (a singleton block repeats a walk
value and leaves the overlap range unchanged --- machine-exact,
\texttt{face\_o5\_gate}/\texttt{face\_b6\_frame}). The mass
evaluation is the prime number theorem with the classical error
term \cite[\S18]{Dav}; the consumption normalization divides it
out. The face verifies both the exact Fubini identity on the live
family and the PNT mass at height ($1.00041$ at $X=2\times
10^{5}$).
\end{proof}

\begin{lemma}[W3: triple-pair factorization]\label{l:W3}
The $\{2,2,2\}$ class aggregate of the sixth-moment ledger
equals the product of the three pair main terms against the joint
overlap volume, up to a deviation aggregate
$\ll_A KJ'V'^2(\log X)^{-A}$ in the consumed normalization.
\end{lemma}

\begin{proof}
The six positions form three disjoint affine pair-chains. The
$p$-local density factorizes over the three chains by the CRT and
Proposition~\ref{p:local} (disjoint slots; the merger
corrections are the pair events already booked in the local law), so the
singular-series model is the product of the three twin series.
Archimedean coupling enters only through the joint window overlap
(the walk geometry $O_6$), which is the certified exact layer of
Proposition~\ref{p:t222}. The deviation from the product of pair
mains is controlled by three applications of the $k=2$ instance
of the spectral transport: by Proposition~\ref{p:prod} the
resolved variance factorizes over the three chains, each factor a
windowed single prime sum, and the three-leaf closure of
\S\ref{s:four} applies to each factor with the exponents of
Proposition~\ref{p:consume}. The face tracks the product model at
$1.9\%$ (median over four shift patterns) at $X=6\times10^{4}$.
\end{proof}

\begin{proposition}[W1: the resolved multi-dimensional
comb]\label{p:W1}
For a $b$-cycle with $b-1$ lock coordinates, the resolved dilated
density
$\rho(\tau_1,\dots,\tau_{b-1})$ obeys the multi-dimensional comb
identification: the model comb is supported on the product CRT
progressions (one congruence per coordinate, as in
Lemma~\ref{l:comb}), with local factors the $\kappa_b$-values of
Proposition~\ref{p:local} and twin-series factors off the moduli;
and the resolved spectral $L^2$-distance from the comb obeys the
bound of Theorem~\ref{t:D} with the same exponents.
\end{proposition}

\begin{proof}
By Proposition~\ref{p:prod} the resolved variance is carried by
the product $\prod_i|S_i|^2$ on the joint analysis grid, and the
multi-$\tau$ transform is complete in each coordinate, so the
$(b-1)$-dimensional Parseval identity holds coordinatewise
(verified exactly in the face, $1.2\times10^{-16}$). The
frequency classification of \S\ref{s:four} applies per factor:
on all-major bins each factor obeys Lemma~\ref{l:MM}; any minor
factor is closed by Lemma~\ref{l:minor}; the comb identification
in each coordinate is Lemma~\ref{l:comb}, whose CRT progression
now runs per coordinate (complete orthogonality in each, since
each transform covers its full period). The argument of
Theorem~\ref{t:D} then repeats verbatim with $b-1$ classification
indices; the truncation exponents strengthen with $k$
(Proposition~\ref{p:consume}). The previously stated bracketing
(aggregate faces above, resolved four-cycle case below) is
subsumed: the aggregate instances are the $\tau$-summed cases,
and the resolved four-cycle case is Lemma~\ref{l:comb} itself.
The resolved-comb face (worked $5$-cycle family) shows the CRT
class contrast directly.
\end{proof}

\subsection{The connected values}\label{s:five:conn}

$C_5=\int O_5C_5^{\mathrm{Urs}}=0.0278(1)$: grouped
midpoint-grid ladder (protocol of \S\ref{s:conv})
$0.02784/0.02781/0.02780$ at $dv=0.05/0.04/0.032$, monotone
settling; the $O$-weighted integration box is exact. The
six-cycle connected value $\{6\}=-0.0078(8)$: midpoint ladder
with the same protocol. (An earlier endpoint-grid evaluation
$C_5=0.0275(5)$, $\{6\}=-0.0100(12)$ is retired: endpoint grids
sign-flip on coarse steps at $b\ge6$ --- register item, 32nd
conversion --- while midpoint grids settle monotonically; the
correction of the four-cycle anchors to their exact rationals,
\S\ref{s:five:cert}, was forced by the same audit.)
Fluctuations are
controlled by the factor-wise variance lemmas (the transports of
Theorem~\ref{t:D}; the resolved multi-dimensional comb exhibits
are bracketed by the aggregate faces and the resolved four-cycle
case --- now Proposition~\ref{p:W1}). End-to-end faces: the
assembled
$m_5=5.6111(1)$ and $m_6=10.1657(18)$ both lie inside the
sine-model bands \eqref{e:guebands}.

\emph{Grading note.} After the certification of
\S\ref{s:five:cert}, the constants still carried numerically are
exactly three: $C_5$, $\{4,2\}$, $\{6\}$ --- evaluations of
explicit piecewise-polynomial integrals in dimensions $4$, $4$,
$5$, at the stated two-sided bands, not certified enclosures.
By the dimension theorem of \S\ref{s:conv} each of the three
is a finite rational combination of rational polytope volumes,
hence \emph{rational}; the exact-integration algorithm of
Proposition~\ref{p:t222} applies to them unchanged, the
computation being heavier by the extra iterated levels and the
ordered-partition count of the connected integrands. Their
closed-form evaluation is the designated next computation
(item N1, \S\ref{s:verif}). The headline of this paper consumes these
constants \emph{with} their bands through the exact dual
certificate of \S\ref{s:lp}, so the claim is exactly: \emph{the
stated bounds hold provided the three connected constants lie
within their stated bands}.

\subsection{The convolution calculus, rationality, and the
midpoint protocol}\label{s:conv}

The connected constants admit a second, structural
representation that both fixes their arithmetic nature and
dictates the numerical protocol. Write $W$ for the unit-square
window autocorrelation and $B$ for the sine-kernel box; the
cumulant integrands assemble from iterated convolutions of these
two profiles. Three results (derivations and verification gates
in the reproduction package, the \texttt{cyclic\_cumulant} and
\texttt{midpoint\_ladder} engines):

\emph{(i) The trace/convolution identity.} At $b=2$ the
connected integral evaluates in closed form,
$\int O_2C_2=1-\|W\star\widehat B\|^{2}=\tfrac13$ exactly ---
the calibration gate (F-TRACE) for everything below; the $b=3,4$
instances reproduce the exact pairing values of
Proposition~\ref{p:t222} to $3\times10^{-14}$.

\emph{(ii) The partition-cyclic cumulant formula (F-CYC).} For
the $b$-cycle connected class,
\[
C_b(v)\;=\;\sum_{P\vdash[b]}(-1)^{|P|-1}
\sum_{\text{cyclic orders of }P}
\operatorname{ov}\bigl(\text{merged-}w\text{ walk}\bigr),
\]
the inner sum running over the distinct cyclic arrangements of
the blocks and $\operatorname{ov}$ the overlap volume of the
merged walk --- validated against the direct Ursell integrands
at $b=3,4,5$ to $3\times10^{-14}$.

\emph{(iii) Dimension and rationality.} In every
$(P,\sigma)$-term of (ii) the walk constraints eliminate one
integration variable ($(m-1)$-dimensionality), and the resulting
domain is a finite union of polytopes cut by the unit-coefficient
hyperplane family $\mathcal H$ of Appendix~\ref{a:cert}(A1)
with polynomial integrands of rational coefficients. Each term
is therefore a rational number, and so are $C_5$, $\{4,2\}$ and
$\{6\}$ --- the same argument that produced
$t_{\mathrm{adj}}=\tfrac7{60}$ and $\{2,2,2\}=\tfrac{131}{420}$
in closed form. What separates the three remaining constants
from Proposition~\ref{p:t222} is only the piece count of the
exact integration, not its kind (item N1, \S\ref{s:verif}).

\emph{The midpoint protocol.} The numerical ladders for
(i)--(iii) revealed a grid pathology recorded as the registry's
$32$nd conversion (Appendix~\ref{a:record}, D10): endpoint
grids on the kinked integrands sign-flip at $b\ge6$ on coarse
steps, while midpoint grids --- which never sample on the kink
set $\mathcal H$ --- settle monotonically. All connected
constants of this paper are evaluated by grouped midpoint-grid
ladders, calibrated at $b=4$ against the exact values of
Proposition~\ref{p:t222} (agreement to the ladder's own band),
and the earlier endpoint-grid adoptions are retired.

\section{Consumption and the proof of
Theorem~\ref{t:main}}\label{s:lp}

With $(m_1,\dots,m_4)=(1,\tfrac43,2,\tfrac{13}4)$ pinned
two-sidedly (Theorem~\ref{t:m3}, Theorem~\ref{t:D},
\eqref{e:moments}), write $C_5=0.0278+\delta$,
$|\delta|\le0.0001$: the fifth- and sixth-moment constraints are
\[
m_5\;\ge\;M_5(\delta):=\tfrac{67}{12}+0.0278+\delta,
\qquad
m_6\;\le\;M_6(\delta):=\tfrac{39}4+\tfrac{131}{420}
+6(0.0278+\delta)-0.0552-0.0078+0.0018,
\]
the same $\delta$ entering both (correlated-band accounting: the
$6C_5$ term of $m_6$ moves with the $C_5$ term of $m_5$; the
independent $m_6$ band $0.0018$ itemizes as $0.0002+0.0008
+0.0008$ over the frozen-slot transport, $\{4,2\}$ and $\{6\}$
--- the former $\{2,2,2\}$ band contribution is retired by
Proposition~\ref{p:t222}). The consumption is no longer a
numerical linear program but an exact certificate:

\begin{lemma}[Exact dual certificate]\label{l:dual}
Let $a=\tfrac{27}{50}$, $b=\tfrac{263}{200}$, $c=\tfrac{103}{50}$
and
\[
P(x)\;=\;\frac{\bigl[(x-a)(x-b)(x-c)\bigr]^{2}}{(abc)^{2}}
\;=\;\sum_{k=0}^{6}y_kx^{k}.
\]
Then $P\ge0$ on $\mathbb R$ (a perfect square), $P(0)=y_0=1$,
$y_5<0$, $y_6>0$. Consequently, for \emph{any} probability
measure $\nu$ on $[0,\infty)$ --- no support bound --- with
$m_1=1$, $m_2=\tfrac43$, $m_3=2$, $m_4=\tfrac{13}4$,
$m_5\ge M_5$, $m_6\le M_6$:
\[
\nu(\{0\})\;\le\;\int P\,d\nu
\;=\;\sum_{k\le4}y_km_k+y_5m_5+y_6m_6
\;\le\;\sum_{k\le4}y_km_k+y_5M_5+y_6M_6\;=:\;w_0(M_5,M_6),
\]
and $w_0$ is affine in $(M_5,M_6)$, hence maximized over the
$\delta$-band at a corner ($\mathrm dw_0/\mathrm d\delta
=y_5+6y_6<0$, so the worst corner is $\delta=-0.0001$:
$M_5=\tfrac{168331}{30000}$, $M_6=\tfrac{42701}{4200}$).
There, exact rational arithmetic gives
\[
w_0\;=\;\tfrac{1153107070889}{11233957316589}
\;=\;0.10264\ldots,
\qquad
1-2w_0\;\ge\;0.7947,\qquad 1-w_0\;\ge\;0.8973 .
\]
\end{lemma}

\begin{proof}
Nonnegativity and $P(0)=1$ are immediate from the closed form;
the sign conditions $y_5=-2(a{+}b{+}c)\cdot(ab{+}bc{+}ca)\ldots$
follow from the expansion with positive roots (displayed in the
certification package; the denominator-cleared expansion, the
corner evaluation, and the final inequalities are
kernel-checked in core Lean~4, \texttt{lean/RhGate/
Certificate.lean}, and verified independently in exact rational
arithmetic by \texttt{certification/certify\_lp.py}).
The mass bound is the Chebyshev--Markov mechanism of
Lemma~\ref{l:cert} at degree six: $y_5<0$ converts the lower
moment bound, $y_6>0$ the upper. The atoms $(a,b,c)$ were
extracted from the numerical LP extremal (a pinched four-atom
measure) and rationalized; any rational triple gives a valid
bound, so the LP solver leaves the trust base.
\end{proof}

Theorem~\ref{t:main} follows through the consumption mechanism of
\S\ref{s:frame}. The rungs: $m_4$ alone gives $13/18$, $31/36$
(Corollary~\ref{c:1318}, certificate $Q$ of Lemma~\ref{l:cert},
leading coefficient positive --- also support-free). The
Dirichlet $L$-function case is verbatim. All constants
ineffective.\qed

\begin{corollary}[Sign-only rung]\label{c:signs}
If the three connected constants merely satisfy $C_5\ge0$,
$\{4,2\}\le0$ and $\{6\}\le0$ (with $C_5\le0.1$), then, by the
re-optimized exact certificate of Appendix~\ref{a:cert}(A5),
\[
\liminf\frac{N_0^{s}}{N}\;\ge\;0.7403,\qquad
\liminf\frac{N_d}{N}\;\ge\;0.8701 .
\]
Sign certification is a far coarser target than an enclosure for
the designated 4D/5D exact integration; the full payoff curve is
tabulated in Appendix~\ref{a:cert}(A5).
\end{corollary}

\begin{remark}[The repaired intermediate rung; register
D6]\label{r:rung}
An earlier stage of this research quoted an intermediate rung
($0.7295/0.8647$) from the fourth- and fifth-moment chain alone,
``consuming no sixth-moment information''. The certification
audit found this rung unsound as stated: with only a \emph{lower} bound on
$m_5$ and no $m_6$ (or support) information, the origin-mass
problem on unbounded support is not improved at all --- mass
$\varepsilon=c/\Lambda^{4}$ at height $\Lambda$ boosts $m_5$
arbitrarily at vanishing cost to $m_1,\dots,m_4$, and the audit's
LP at support ranges $8/20/60/200$ degraded
$0.73144\to0.72580\to0.72340\to0.72257\to\tfrac{13}{18}$. Any
dual certificate for a lower fifth-moment bound needs a negative
$x^5$-coefficient dominated by a positive $x^6$-coefficient, i.e.\
sixth-moment (or support) information is \emph{necessary}. The
rung is therefore demoted: it holds in the stated form only with
the $m_6$ cap, in which case it is subsumed by
Lemma~\ref{l:dual}. The full certificate is unaffected
($y_6>0$).
\end{remark}

\medskip
\noindent\emph{Merge-audit note on the band model.} The audit
first ran an \emph{anti}-correlated variant (the $\delta$ of
$m_5$ entering $m_6$ with opposite sign) and obtained $0.78972$;
the correlated form above is the correct accounting, since $C_5$
is a single unknown shifting both constraints together. Both runs
are recorded in the register (D3) so the sensitivity of the
headline to the band model is on the record: under the
(incorrect, doubly conservative) anti-correlated accounting the
headline would read $0.7897/0.8949$.

\section{Numerical verification}\label{s:faces}

The standalone pipeline (\texttt{pipeline/}; \texttt{numpy},
\texttt{scipy}) reruns every gate in a few minutes. Central
height table (sixteen coprime modulus families):

\begin{center}
\begin{tabular}{lcccccccc}
$T$ & 2400 & 4800 & 9600 & 19200 & 38400 & 57600 & 76800 & 153600\\
\hline
dispersion ($10^{-3}$) & 8.39 & 5.19 & 4.16 & 4.06 & 3.87 & 1.84
& 1.29 & 1.04\\
tracking (\%) & 1.5 & 1.0 & 0.7 & 0.9 & 0.8 & 0.4 & 0.2 & 0.2\\
\end{tabular}
\end{center}

\noindent (full-range geometric mean $0.706$ per doubling: the
pure averaging law $2^{-1/2}$). Further faces: local closure
exact (zero violations, seven configurations); Parseval exact;
per-bin Siegel--Walfisz face falling tenfold over the range;
Bell(5)/Bell(6) gates exact; five- and six-fold product faces
green at all heights; well-spacing zero violations; hybrid size
control $24$--$31\times$ height-stable; consumption LP
grid-stable. The registry: $194$ pre-registered checks, $160$
passed, $34$ fired and converted, all forward-recorded (the
$31$st--$34$th conversions belong to the corrected-constants
pass and the Lean compiles, \S\ref{s:verif}(f),
Appendix~\ref{a:record}).

The certification audit of August 16, 2026 re-ran: the full
moment-chain gate suite at $T=9600$ (ALL PASS, $72$s: dispersion $4.163\times
10^{-3}$, tracking median $0.73\%$, SW RMS $1.599\times10^{-3}$,
tightness $24.8$--$31.3\times$, five-fold $0.06$--$1.29\%$,
six-fold $0.01$--$0.73\%$, LP $0.79454/0.89727$ at two grids
under the since-retired endpoint constants; the
corrected-constants rerun gives $0.79472/0.89736$);
the fourth-moment identity suite at $T=9600$ (ALL PASS,
$48$s); the constant suite of \S\ref{s:fm}
(\texttt{m1\_suite.py validate}: $6/6$ local
factors, $100/100$ vector sieve, anchors and cross-convention
difference reproduced; core $-0.00545$ against archive
$-0.00544$; $\alpha$-truncation $-0.00491$); the sawtooth
certification (\texttt{sawtooth\_cert}: grid $1.8675$, samples
$1.8419/1.8494/1.8512$, certificate $2.100$); the continuum
quadrature gates (\texttt{quadrature\_cert}: geometric tail ratio $0.500$
certified, $c_0=-1.4228$ decade-stable); and the certificate
suite of \S\ref{s:cert}
(\texttt{certificate\_verification.py} exact;
\texttt{moment\_lp\_reopt} $0.0721$; \texttt{mu3\_ledger\_gates};
\texttt{cumulant\_engine} gates $C_2$ exact,
$\int O_3C_3=0$). An independent re-implementation of the
consumption LP (different code, HiGHS backend) reproduced every
rung as recorded in \S\ref{s:lp}.

A second certification pass (same date) added: the exact
pairing-layer integrals (\texttt{exact\_t222.py}: all values of
Proposition~\ref{p:t222}, dual-scheme identical rationals, the
closed-form cross-checks, and the match of the conjectured
rationals); the exact dual certificate
(\texttt{certify\_lp.py}: expansion, both band corners, final
inequalities, all exact); the write-out faces
(\texttt{writeouts\_verification.py}: free-slot Fubini exact,
window mass $1.00041$, triple-pair product model median $0.9814$,
resolved Parseval $1.2\times10^{-16}$, resolved CRT contrast ---
one face fired during construction and was converted, the
registry's $30$th); and the support-range probe behind
Remark~\ref{r:rung}.

The corrected-constants pass re-ran: the grouped
midpoint-grid ladders for $C_5$, $\{4,2\}$, $\{6\}$
(\S\ref{s:conv}; checkpointed engines, $b=4$ calibration exact
against Proposition~\ref{p:t222}); the exact dual certificate at
the corrected corners (\texttt{certify\_lp.py}: expansion,
both band corners, final inequalities, all exact, worst corner
$w_0=\tfrac{1153107070889}{11233957316589}$); the convolution
trace gates (F-TRACE: $b=2$ exact $\tfrac13$, $b=3,4$ to
$3\times10^{-14}$); and --- for the first time in the program
--- a genuine Lean kernel compile of the identity layer on the
maintainers' toolchain (\S\ref{s:verif}(f)).

\begin{remark}[Effective faces, ineffective theorem]
The faces are finite-height measurements with effective
constants; the theorem's constants are ineffective (Siegel). The
faces verify the exact-identity layer, the model identifications,
and the qualitative averaging behaviour the proof predicts; they
cannot and do not verify the asymptotics. The two kinds of
evidence are kept distinct throughout.
\end{remark}

\section{Verification status}\label{s:verif}

\emph{Claim grade: certified-candidate.} The chain = classical
inputs (Siegel--Walfisz \cite{Dav,IK}, Vaughan, singular-series
truncation, \cite{MRT} for \S\ref{s:fm}) + archived certified
lemmas (bridge, glue, aggregation, D1 truncation, consumer
interface, anchor certificates: \S\ref{s:fm}) + the
machine-verified identities of
\S\S\ref{s:cert},~\ref{s:local}--\ref{s:five} + the numerically
evaluated connected constants of \S\ref{s:five:conn} + the
consumption layer of \S\ref{s:frame}. Grades by layer:

\emph{(a) Machine-exact or exact-rational, re-verified in the
certification audit.} The certificate algebra (Lemma~\ref{l:cert}) and
$\kappa(c)$ laws; the local closure
(Proposition~\ref{p:local}); the Parseval and product identities;
the Bell ledgers and frozen-slot identities; the anchor-free
$67/12$ and $39/4$; the $197/60$ ledger arithmetic and the now
fully exact two-bookkeeping reconciliation
(Remark~\ref{r:books}); the exact pairing-layer constants
$t_{\mathrm{adj}}=\tfrac7{60}$, $t_{\mathrm{opp}}=\tfrac1{30}$,
$\{2,2,2\}=\tfrac{131}{420}$ (Proposition~\ref{p:t222},
Appendix~\ref{a:cert}); the exact dual certificate
(Lemma~\ref{l:dual}) --- the consumption step involves no LP
solver, no grid, and no support bound; the sawtooth constant
$C_{\mathrm{saw}}=2.10$ (grid-plus-analytic certification).

\emph{(b) Certified-candidate analytic layer (pre-Lean,
pre-review).} The nine-class $\mu_3$ ledger (two nonstandard
steps written out in the archived memo); the truncation
discharge; the spectral proof of Lemma D (Theorem~\ref{t:D}) ---
whose minor/mixed-bin size bookkeeping is classical in kind but
carried at write-out grade with face-checked budgets --- and its
factor-wise transports to $k=5,6$, including the written-out
W1--W3 (\S\ref{s:writeouts}), which are proved at this layer's
grade; the covered-zone/band/glue chain of \S\ref{s:fm}. Five
adversarial audit rounds precede the frozen chain (archive
rounds r93, r95, r100, r102).

\emph{(c) Write-out items.} Discharged:
W1 (Proposition~\ref{p:W1}), W2 (Lemma~\ref{l:W2}), W3
(Lemma~\ref{l:W3}), each with a pre-registered face. W2 is
complete (elementary); W1 and W3 are proved at the grade of
layer (b), since they consume the same spectral transports.

\emph{(d) Numerical constants, not yet certified enclosures.}
N1 (reduced, and now known rational): $C_5$, $\{4,2\}$,
$\{6\}$ (grouped midpoint-grid ladders with stated bands,
\S\ref{s:five:conn}; each proved a rational number by the
dimension theorem of \S\ref{s:conv}, so certification means
identifying a fraction, not enclosing a real); the designated
route is the exact-integration algorithm of
Proposition~\ref{p:t222} (same kink structure; heavier by
dimension and by the ordered-partition count --- cluster scale,
mechanical in kind). The \emph{certificate family} of
Appendix~\ref{a:cert}(A5) makes the payoff graded and
plug-in: any future enclosure yields an exact certificate, and a
sign-only certification already yields $0.740/0.870$. N2
(partially discharged): the quadrature's slope
is proved exactly $1$, its constant is identified in closed form,
$c_0=\gamma-2$ (Proposition~\ref{p:c0}), and the geometric tail
was already certified (r60); what remains computed is the mean
oscillation $\overline{\mathrm{osc}}$ inside $c^{*}$ and the
universality-collapse rate write-out. Only the conservative
charge $\varepsilon_{CL}$ is consumed anywhere
(Lemma~\ref{l:CL}). Full certification of N1 upgrades
Theorem~\ref{t:main} to the grade of Corollary~\ref{c:1318};
of N2, Theorem~\ref{t:fm} to $0.7031/0.8515$.

\emph{(e) Imports.} \cite{C26}: the framework of
\S\ref{s:frame} (matrix, block structure, tail, counting,
first two moments) --- a publicly posted preprint of this
program, not independently peer-reviewed; its own analytic
imports are published literature. Directly consumed published
literature: \cite{Wei52} and \cite[Thm.~5.12]{IK} (explicit
formula), \cite[Thm.~9.2]{Tit86} (zero density),
\cite{Mon73,BGST} (prime-side pair-correlation evaluation,
unconditional in the band), \cite{MRT}, \cite{HR},
\cite{Dav}.

\emph{(f) Formalization status.}
The identity layer has now been \emph{compiled and
kernel-checked} on the maintainers' machine: core Lean~4
(toolchain \texttt{leanprover/lean4:v4.33.0}, no
\texttt{mathlib}), package \texttt{lean/RhGate/}, build log
included. Compiled content: the anchor-free assembly identities
$m_4=\tfrac{13}4$, $m_5$-part $=\tfrac{67}{12}$, $m_6$
pair-layer $=\tfrac{39}4$ as universally quantified rational
identities (any anchors $t_{\mathrm a},t_{\mathrm o}$ ---
\texttt{Anchors.lean}); the consumption arithmetic
$\tfrac{13}{18}$, $\tfrac{31}{36}$; and the coincidence-counting
local law at $b=4$, $p=5$ by full kernel enumeration of all
$125$ lock classes (\texttt{LocalClosure.lean} --- this theorem
is \emph{axiom-free}; the others use only
\texttt{propext}, \texttt{Classical.choice},
\texttt{Quot.sound}). Zero \texttt{sorry}s. The compile itself
fired a registered check and became the $33$rd conversion: the
seed files' proofs invoked \texttt{ring} and \texttt{decide}
on $\mathbb Q$, neither available in core Lean (\texttt{ring}
is a \texttt{mathlib} tactic; core \texttt{Rat} arithmetic is
not kernel-reducible, so \texttt{decide} stalls on
\texttt{instDecidableEqRat}); all five affected proofs were
repaired with the core \texttt{grind} tactic, statements
unchanged (Appendix~\ref{a:record}). A second same-machine compile (August 16, 20:09) covered the
certificate layer: \texttt{Certificate.lean} (the
denominator-cleared expansion identity of Lemma~\ref{l:dual},
nonnegativity, normalization, both corner evaluations of $w_0$,
the corner-selection inequality $y_5+6y_6<0$ --- so the
\emph{choice} of binding corner is itself kernel-checked ---
the headline inequalities $1-2w_0\ge\tfrac{7947}{10000}$,
$1-w_0\ge\tfrac{8973}{10000}$, the monic expansion of
Appendix~\ref{a:cert}(A6), the exact $\{2,2,2\}$ arithmetic
and the instantiated anchor identities) and
\texttt{LocalClosure2.lean} (the local law at $b=4$, $p=7$ and
$b=5$, $p=5$: $343$ and $625$ lock classes by kernel
\texttt{decide} on $\mathbb N$ --- both \emph{axiom-free}).
This compile fired the registry a second time (the $34$th
conversion): \texttt{grind}'s ring/linear core treats a square
$(\cdot)^2$ as an opaque atom and cannot close perfect-square
nonnegativity; the repair is the \emph{core} ordered-ring lemma
\texttt{Lean.Grind.OrderedRing.sq\_nonneg} (term-mode,
statement unchanged, still mathlib-free). The pre-compile seed
\texttt{SOSCertificate.lean} --- whose headline constants
carried the retired endpoint-grid corner --- is removed from the
package to preclude coexisting $0.7946$/$0.7947$ headline
theorems; its still-valid content is absorbed into
\texttt{Certificate.lean}, and the draft is preserved in the
program archive. Every statement is proxy-verified in
exact rational arithmetic (\texttt{certify\_lp.py}) or full
\texttt{Python} enumeration (zero deviations). What Lean
certifies is the rational-arithmetic layer alone: the moment
pinning, the measure-theoretic Chebyshev--Markov step and the
analytic chain remain graded as in layers (b)--(e). The truncation, spectral and consumption
\emph{analysis} remains unformalized; formalization of the
analytic layer --- which relies on the framework of
\cite{C26} --- is graded pending. The formalization package
accompanying \cite{C26} is separate from this repository; the
compiled coverage of this paper is exactly as stated above.

\emph{(g) The gate.} Per the program's verification gate, the
record claim follows (i) the certified fractions N1 (and
optionally N2), (ii) compiled Lean formalization of the identity
and certificate layers with the queue's Q1--Q6 (both are now
compiled, \S\ref{s:verif}(f); the analytic layer remains), and
(iii)
external review. No Riemann Hypothesis, zero-density input,
pair-correlation conjecture, or Kloosterman technology is used
anywhere. Nothing here bears on RH itself.

\section*{Acknowledgements}

The mathematical development of this paper --- the framework of
\S\ref{s:frame}, the certificate calculus, the spectral proof of
Lemma D, the moment identifications, the convolution calculus,
and the reproduction and Lean packages --- was carried out by
Claude (Anthropic), a large language model, working within a
research program directed by the authors under the
falsifier-first verification discipline described in
\S\ref{s:verif} and Appendix~\ref{a:record}. The authors
specified objectives, arbitrated decisions, compiled and audited
the Lean package on their own toolchain, and take full academic
responsibility for the content. This paper is built on the
two-thirds theorem \cite{C26}, a preprint of the same research
program.

\appendix

\section{The comb identification, with a worked
pair}\label{a:comb}

In the $\tau$-frame the model of $\rho(\tau)$ is assembled from
twin singular series: with $\SSS$ the Hardy--Littlewood series
$\SSS(h)=2C_2\prod_{p\mid h,\,p>2}\frac{p-1}{p-2}$ ($h$ even),
window overlap counts $V_m,V_n$, and exact diagonal pieces,
\[
\mathrm{Model}(\tau)=\!\!\sum_{\substack{t\equiv0\ (b_2)\\
t\equiv\tau\ (b_1)\\ t\notin\{0,\tau\}}}\!\!
\SSS\!\Bigl(\frac{t}{b_2}\Bigr)\SSS\!\Bigl(\frac{t-\tau}{b_1}\Bigr)
V_m\!\Bigl(\frac{t}{b_2}\Bigr)V_n\!\Bigl(\frac{t-\tau}{b_1}\Bigr)
+\ \text{diagonals}.
\]
The CRT progression $t\equiv0\ (b_2)$, $t\equiv\tau\ (b_1)$ is
the \emph{complete} orthogonality of the joint frame: a shift $t$
contributes iff it is simultaneously a $b_2$-lattice twin shift
on the $m$-side and, offset by $\tau$, a $b_1$-lattice twin shift
on the $n$-side; no incomplete sums occur because the
$\tau$-transform runs over the full period. The local factors of
$\SSS$ at the primes $p\mid b_1b_2$ are exactly the
$\kappa$-values of Proposition~\ref{p:local} at $b=4$; for
$p\nmid b_1b_2$ they are the twin factors.

\emph{Worked pair} $(b_1,b_2)=(5,3)$: $M=15$, and for
$\tau\equiv\tau_0$ the surviving shifts are the single residue
class $t\equiv t_0(\tau)\ (15)$ with
$t_0=3\cdot(2\tau\bmod5)\bmod15$ (here $3^{-1}\equiv2\ (5)$). For
example $\tau=1$ gives $t\equiv6\ (15)$: the model sums
$\SSS(t/3)\SSS((t-1)/5)V_mV_n$ over $t=6,21,36,\dots$ and the
negative branch, plus the diagonal $t=\tau$ term when $3\mid\tau$
and $t=0$ when $5\mid\tau$ (neither here). This is precisely what
\texttt{face\_dispersion} constructs; its agreement with the
measured profile at $0.2$--$4\%$ across all sixteen families and
eight heights is the numerical face of Lemma~\ref{l:comb}. At
cycle lengths $5$ and $6$ the same complete-orthogonality logic
gives the product-model combs of \S\ref{s:five} (aggregate faces
green at $0.1$--$2.5\%$).

\section{Provenance, research record, and discrepancy
register}\label{a:record}

\subsection*{The research record}

The mathematics of this paper was developed in an interactive
research program (August 2026) whose working discipline is
falsifier-first: every structural claim was registered with a
numerical check before adoption, every firing was converted on
the forward record, and the complete trail --- round logs,
memos, claim ledgers, correction history, scripts --- is
preserved in the program archive that accompanies the
reproduction package. The four stages of the chain
(\S\ref{s:cert}, \S\ref{s:fm},
\S\S\ref{s:trunc}--\ref{s:four},
\S\S\ref{s:five}--\ref{s:lp}) were each closed under
adversarial audit rounds before the certification audit of
\S\ref{s:verif} froze the chain presented here.

\subsection*{Conversions that shaped this text}

Six of the registry's fired-and-converted checks shaped the
mathematics directly: the evaluation-transfer repair (envelopes
are never transported to limits; r61); the model-subtraction
correction of the reduced covariance (r80); the hybrid size
budget (two firings, r88); the $\tau$-frame lesson (fixed-cutoff
arc statistics never self-average; the Parseval identity entered
here; r90); the incomplete-sum slip of the progression frame
(r95, Remark~\ref{rem:prog}); and the removal of an unproven
sixth-moment cap from an internal ladder rung (the intermediate
values quoted in Theorem~\ref{t:main} consume exactly what their
chains prove). The 29th conversion (r103) is the Lean-draft
inverse omission of \S\ref{s:local}.

\subsection*{Discrepancy register of the certification audit
(August 16, 2026)}

\begin{itemize}
\item[D1.] Earlier draft of \S\ref{s:fm}, assembly paragraph:
printed $\Lambda_2(0)=0.156293$; the closed form at $\bar R=0.0066$ gives
$0.154193$, and the theorem constants $0.6916/0.8458$ are
consistent with the latter. Corrected in \S\ref{s:fm}; theorem
unaffected. \emph{Impact on the main theorem: none.}
\item[D2.] An earlier draft attributed the bands
$5.573\pm0.113$, $10.02\pm0.27$ to measurement ``on
$2\times10^{7}$ zeta zeros''; the reproduction package derives
exactly these numbers from GUE double-window Richardson
extrapolation (\texttt{gue\_bands.py},
\texttt{m5m6\_bands.py}), and no zeta-zero
dataset of that size exists in the record. This paper
attributes the bands to the sine-model/GUE measurement
\eqref{e:guebands}. If a zeta-zero measurement exists it is not
in the record and is not consumed. \emph{Impact: attribution
only; no constant changes.}
\item[D3.] Consumption LP band model: correlated vs
anti-correlated $\delta$-accounting ($0.79454$ vs $0.78972$);
the correlated form is correct and is what
\texttt{face\_lp\_full} implements; both on record
(\S\ref{s:lp}). \emph{Impact: the anti-correlated variant is
doubly conservative; headline unchanged.}
\item[D4.] The $m_4$+$m_5$ rung: an earlier stage quoted
$0.7295/0.8647$; the audit's independent LP gives
$0.73144/0.86572$ under the same constraints at bounded support.
Superseded by D6: the rung is demoted regardless.
\emph{Impact: none on the main chain (the demoted rung is not
consumed).}
\item[D5.] Formalization language: earlier program documents
described \cite{C26} as ``Lean-formalized''; that statement
concerns the package accompanying \cite{C26}, of which no
compiled artifact exists in this repository
(\S\ref{s:verif}(f)). The compiled coverage of this paper is
exactly as graded in (f).
\emph{Impact: grading language only.}
\item[D6.] The intermediate $m_4$+$m_5$ rung
($0.7295/0.8647$, an earlier stage) is unsound on unbounded
support without sixth-moment information: the audit's LP degrades
to $13/18$ as the support range grows ($0.73144\to0.72257$ at
$R=8\to200$), and a dual certificate for a fifth-moment lower
bound necessarily consumes $x^6$-information. Rung demoted and
repaired (Remark~\ref{r:rung}); the main chain is unaffected
($y_6>0$ in Lemma~\ref{l:dual}). \emph{Impact: removes an
intermediate rung; headline unaffected.}
\item[D7.] The archived four-cycle anchor
$t_{\mathrm{adj}}=0.11717$ (rounds 69--71; quoted in earlier
drafts and \texttt{m1\_suite}) carries a $+0.43\%$
rectangle-rule grid bias: the exact value is $\tfrac7{60}
=0.11667$ (Proposition~\ref{p:t222}). Downstream totals are
unaffected (the $m_5$/$m_6$ rational layers are anchor-free);
the saturated-split triple $0.2677/0.0156/{-0.0177}$ becomes
exactly $\tfrac4{15}/\tfrac1{60}/{-\tfrac1{60}}$
(Remark~\ref{r:books}), in better agreement with the cumulant
engine's independent $-0.01663$. \emph{Impact: none on the
rational layers (anchor-free); corrects the archived anchor.}
\item[D8.] The rational identifications
$T_0,\dots,T_3=\tfrac3{70},\tfrac1{90},\tfrac1{180},\tfrac1{70}$,
recorded in the audit ledger as a conjecture and not consumed, are
now proven (Proposition~\ref{p:t222}); $\{2,2,2\}=\tfrac{131}{420}$
is consumed exactly and its band contribution retired.
\emph{Impact: tightens $m_6$; contributes to the headline.}
\item[D9.] The continuum-quadrature constant
$c_0=-1.4228$ (r60, ``computed, not fitted'') is identified in
closed form: $c_0=\gamma-2$ (Proposition~\ref{p:c0}). The audit
located a $d=2$ discrepancy between the archive's tensor-sieve
$\gamma$-array and the raw Euler closed form
($\widetilde\gamma_2=C_2-1$ vs $\gamma_2=C_2$, shifting the
constant by exactly $-2$); this is the archive's parity-mean
convention (the $q=2$ arc block booked in $\bar D$), not an
error --- the raw-convention constant is $\gamma$, and both
conventions are now documented with the conversion. The slope-$1$
law of the partial sums, previously a measurement
($0.99996$), is proved exactly. \emph{Impact: N2 partial
discharge; main theorem unaffected.}
\item[D10.] The endpoint-grid ladders behind
the previous connected constants fired their refinement face:
at $b\ge6$ and coarse steps the endpoint evaluation
\emph{sign-flips} ($b=6$ at $dv=0.4$: $+0.036$ against the true
$-0.010$), and the archived adoptions $C_5=0.0275(5)$,
$\{4,2\}=-0.0548(10)$, $\{6\}=-0.0100(12)$ inherited a
systematic shift. Grouped \emph{midpoint}-grid ladders
(\S\ref{s:conv}), calibrated exactly against
Proposition~\ref{p:t222} at $b=4$, settle monotonically and give
the corrected $C_5=0.0278(1)$, $\{4,2\}=-0.0552(8)$,
$\{6\}=-0.0078(8)$ consumed in \S\ref{s:lp} (registry's
$32$nd conversion; the $31$st is D7's anchor correction,
independently forced by the same audit). \emph{Impact: direct
--- these are the constants the headline consumes; the theorem
value moved $0.7946\to0.7947$ under the correction.}
\item[D11.] The Lean draft files claimed
core-decidability of their $\mathbb Q$-statements
(``\texttt{decide} on $\mathbb Q$'') and used \texttt{ring};
the maintainers' compile showed both unavailable in core
Lean~4.33 (\texttt{ring} is \texttt{mathlib};
\texttt{instDecidableEqRat} does not kernel-reduce). Five
proofs repaired with \texttt{grind}, statements unchanged;
build then succeeded (registry's $33$rd conversion;
\S\ref{s:verif}(f)). \emph{Impact: proofs only; statements
unchanged.}
\item[D12.] The second compile fired on
\texttt{cert\_nonneg}: \texttt{grind} assigns the square
$(\cdot)^2$ as an opaque atom (its diagnostics guess the atom's
value) and cannot derive perfect-square nonnegativity; repaired
with core \texttt{Lean.Grind.OrderedRing.sq\_nonneg}
(registry's $34$th conversion). The audit also found the
superseded seed \texttt{SOSCertificate.lean} silently outside
the build roots while declaring identically named headline
theorems at the retired constants; the file is removed from the
package (content absorbed; the draft is preserved in the
program archive) and
the build now covers every \texttt{.lean} file in the
directory. \emph{Impact: proofs and packaging only; statements
unchanged.}
\item[D13.] The citation relation to \cite{C26} was reversed
twice on the forward record: an intermediate revision absorbed
the framework into the text with proofs (when \cite{C26} was
not yet publicly posted), and the present revision restores
\cite{C26} as the cited foundation and removes the duplicated
proofs, keeping statements only (\S\ref{s:frame}).
\emph{Impact: attribution and exposition; no mathematical
content changed.}
\end{itemize}

\noindent The registry's $30$th conversion occurred during the
certification audit: the first version of the W1 resolved-comb
face
used raw multiples of $b_1b_2$ as the on-progression and fired;
the corrected face implements the per-coordinate CRT classes
(\S\ref{s:writeouts}). Forward-recorded in
\texttt{writeouts\_verification.py}. The $31$st--$34$th are
D7, D10, D11 and D12 above.

\subsection*{Reproduction}

The complete package --- paper, reproduction pipelines,
certification layer, GPU engines, and the Lean sources with their
build log --- is public at
\url{https://github.com/JoshuaHKU/zeta-0.7947-reproduction}.
A consolidated reproduction checklist accompanies this paper
(\texttt{REPRODUCTION.md}): environment, one command per suite,
expected outputs with the August 16, 2026 rerun values, and the
mapping from every quoted constant to the script that regenerates
it. Total runtime of the full set is under fifteen minutes on one
CPU (the session package's optional long gates under thirty; the
exact $\{2,2,2\}$ certification under two hours).

\section{The exact certification layer}\label{a:cert}

\subsection*{A1. The exact-integration algorithm}

Every pairing-layer constant is an integral
$\int_{[-1,1]^{d}}(1-\mathrm{spread}(p_1,\dots,p_b))_+\,
\prod_i|x_i|\,dx$ where the walk positions $p_j$ are
$\{0,1\}$-combinations of the chord variables and $0$ is a
position (so the support forces $|x_i|\le1$ and the
$\min(|\cdot|,1)$ caps of the $C_2$ factors never bind). The
integrand is piecewise polynomial; its kink set is contained in
the hyperplane family
\[
\mathcal H=\bigl\{p_i-p_j=c\bigr\}_{i<j,\ c\in\{0,\pm1\}}
\cup\bigl\{x_i=0,\pm1\bigr\},
\]
all with coefficients in $\{-1,0,1\}$ and unit leading
coefficients in the innermost variable. Iterated integration:
the inner integral is computed exactly between the enumerated
inner breakpoints; because leading coefficients are units, the
breakpoints are themselves small-integer linear forms in the
outer variables, and the piece structure of a partial integral
changes only on (i) direct kinks, (ii) collisions of two inner
breakpoints, (iii) breakpoint--boundary collisions --- an
explicitly enumerated resultant family at each level. On every
kink-free piece the (partial) integrand is a polynomial of
degree at most $2/4/8$ in the inner/middle/outer variable,
integrated exactly by closed Newton--Cotes rules of order
$5/7/9$ at rational nodes, in exact rational arithmetic
throughout. \emph{Self-check}: each piece is integrated by the
same rule on the whole piece and on its two halves; for a
polynomial of degree below the rule order both are the exact
integral, so the two rationals must be \emph{identical}, and a
kink missed by the enumeration forces a discrepancy on the
containing piece. No discrepancy occurred at any level of any
run. Independent cross-checks: closed-form orthant evaluations
(A2), and the float grid ladders of the archive
($10^{-6}$-agreement with the certified rationals).

\subsection*{A2. Closed-form cross-checks}

$T_0$ (positions $\{0,v,w,u\}$): in the all-positive orthant
$\int_{[0,1]^3}(1-\max(v,w,u))vwu=\tfrac1{56}$ (order and
integrate); a one-negative orthant gives
$\int_0^1t(1-t)^5/20\,dt=\tfrac1{840}$; summing the orbit,
$T_0=2\cdot\tfrac1{56}+6\cdot\tfrac1{840}=\tfrac3{70}$.
$t_{\mathrm{opp}}$ (positions $\{0,v,v{+}w,w\}$): the same
decomposition gives $\tfrac1{30}$ exactly, confirming the
archived exact anchor. $t_{\mathrm{adj}}$: the two same-sign
quadrants give $\tfrac1{20}$ each and the two mixed quadrants
$\tfrac1{120}$ each: $t_{\mathrm{adj}}=\tfrac{14}{120}
=\tfrac7{60}$ (not the archived grid value $0.11717$; register
D7).

\subsection*{A3. The certified values}

\[
t_{\mathrm{adj}}=\tfrac7{60},\quad
t_{\mathrm{opp}}=\tfrac1{30},\quad
T_0=\tfrac3{70},\quad
T_1=\tfrac1{90},\quad
T_2=\tfrac1{180},\quad
T_3=\tfrac1{70},\quad
\{2,2,2\}=\tfrac{131}{420},
\]
with $\Phi_4=-\tfrac1{60}$, saturated $\{2,2\}=\tfrac4{15}$,
plateau $=\tfrac1{60}$, connected model $=-\tfrac1{60}$
(Remark~\ref{r:books}). Code:
\texttt{certification/exact\_t222.py} (classes run separately;
dual-scheme identical-rational check at every piece).

\subsection*{A4. The continuum-quadrature constant in closed form
(N2, partial discharge)}

\begin{proposition}[$c_0=\gamma-2$]\label{p:c0}
Let $\gamma_d$ denote the free coefficients of the partial-sum
identity (Lemma~\ref{l:PS}) at the $(1,1)$ class. Then
$\gamma_d$ has the exact Euler closed form
\[
\gamma_{2m}=C_2\prod_{p\mid m}\frac1{p(p-2)}
\quad(m\ \text{odd squarefree}),\qquad
\gamma_d=0\ \text{otherwise},
\]
and, in the raw convention,
$\sum_{d\le M}d\gamma_d=\log M+\gamma+o(1)$; in the archive
convention (the $q=2$ arc block booked inside the parity mean,
$\widetilde\gamma_2=\gamma_2-\beta_2=C_2-1$),
\[
c_0\;=\;\gamma-2\;=\;-1.4227843\ldots
\]
\end{proposition}

\begin{proof}
(i)~The closed form: from the arc identity
(Proposition~\ref{p:S}) at $b_1=b_2=1$,
$\beta_q=\mu(q)^2/\varphi(q)^2$, and
$\gamma_d=\sum_e\mu(e)\beta_{de}$ factorizes over primes with
local weights $f_0(p)=1-(p-1)^{-2}$, $f_1(p)=(p-1)^{-2}$,
$f_{v\ge2}=0$; $f_0(2)=0$ forces $2\,\|\,d$ and the product
assembles to the display ($C_2$ the twin-prime constant).
(ii)~The constant: $\sum_{d\le M}d\gamma_d
=2C_2\sum_{m\le M/2}\prod_{p\mid m}\tfrac1{p-2}$, whose Dirichlet
series is $\zeta(s+1)E(s)$ with
$E(0)=\prod_p(1-\tfrac1p)\prod_{p>2}(1+\tfrac1{p-2})
=\tfrac1{2C_2}$ (slope $2C_2E(0)=1$ exactly) and
$E'(0)/E(0)=\sum_p\tfrac{\log p}{p-1}
-\sum_{p>2}\tfrac{\log p}{p-1}=\log2$ (the $p$-sums cancel
pairwise except $p=2$); the double-pole residue gives
$2C_2[E(0)(\log\tfrac M2+\gamma)+E'(0)]=\log M+\gamma$.
(iii)~The convention shift: the archive tensor sieve
(\texttt{tail\_bound}) and the closed form agree at \emph{every}
$d$ except $d=2$, where the archived class-$d$ replacement books
$\widetilde\gamma_2=C_2-1$ (the $q=2$ block moved to the parity
mean); this shifts $\sum d\gamma_d$ by exactly $-2$. Numerical
verification (\texttt{n2\_c0\_certification.py}): the raw sum
converges to $\gamma$ with deviations
$9.4\times10^{-4}/6.3\times10^{-5}/9.2\times10^{-6}$ at
$M=10^4/10^5/2\times10^7$; the archived measurement $-1.42279$
at $M=1.6\times10^{7}$ matches $\gamma-2$ to $10^{-5}$; the
partial-sum identity closes with its $d>M$ tail
($\mathrm{LHS}-\mathrm{RHS}-\mathrm{tail}=-0.010$ at
$M=2\times10^5$, $D=2\times10^{7}$, consistent with the
remaining truncation).
\end{proof}

\noindent Consequently the ``computed constants'' of
Lemma~\ref{l:CL} reduce to: slope $=1$ (proved), $c_0=\gamma-2$
(closed form), the certified geometric tail $T_\Gamma$
(archive r60, ratio $0.500$), and the single remaining computed
piece $\overline{\mathrm{osc}}$ inside $c^{*}$. N2 is thereby
partially discharged; what remains is the
$\overline{\mathrm{osc}}$ identification (or certified
enclosure) and the universality-collapse rate write-out.

\subsection*{A5. The N1 certificate family and its payoff curve}

The exact dual certificate re-optimizes against any hypothetical
certified enclosures $C_5\in[c_-,c_+]$, $\{4,2\}\le u_{42}$,
$\{6\}\le u_{6}$ (correlated accounting: the single unknown
$C_5$ enters $m_5$ and $m_6$ together; the worst case is a
corner of the affine family). Exact rational certificates
(\texttt{certificate\_family.py}, atoms re-optimized per
scenario):

\begin{center}
\begin{tabular}{lcc}
scenario (what the 4D/5D exact integration certifies) &
$1-2w_0\ge$ & $1-w_0\ge$\\
\hline
current bands ($\pm0.0001$; $u_{42},u_6$ at band tops) &
$0.79472$ & $0.89736$\\
enclosure width $0.01$ & $0.77418$ & $0.88709$\\
enclosure width $0.04$ & $0.74299$ & $0.87149$\\
signs only: $C_5\ge0$, $\{4,2\}\le0$, $\{6\}\le0$ &
$0.74031$ & $0.87016$\\
no connected information & $13/18$ & $31/36$\\
\end{tabular}
\end{center}

\noindent Sensitivity of the shipped certificate (exact):
$\partial(\text{headline})/\partial c_-=+7.32$,
$\partial/\partial c_+=-5.61$,
$\partial/\partial(u_{42}{+}u_6)=-0.93$. In particular a
\emph{sign-only} certification of the three connected constants
--- far coarser than an enclosure --- already certifies
$0.740/0.870$, and the designated 4D/5D exact computation has a
graded payoff at every achievable width. Computational spec from
the certified 3D runs: the exact integrator's cost is
(pieces$\times$nodes) per level ($\approx40$--$400$ per level in
3D at $\le90$ kink forms); the $C_5$ and $\{4,2\}$ integrands
carry $\le541$ resp.\ $\le75\times$(spectator) overlap terms on
the same $\{0,\pm1\}$-kink family, one extra level ($\{6\}$:
two), placing the runs at cluster scale rather than laptop
scale.

\subsection*{A6. The exact dual certificate}

The polynomial of Lemma~\ref{l:dual} expands (numerator, monic)
as
\[
x^6-\tfrac{783}{100}x^5+\tfrac{975601}{40000}x^4
-\tfrac{19203237}{500000}x^3+\tfrac{1599367847}{50000000}x^2
-\tfrac{16571397771}{1250000000}x
+\tfrac{534950348409}{250000000000},
\]
with $(abc)^2=\tfrac{534950348409}{250000000000}$, so
$y_0=1$, $y_5=-\tfrac{783}{100}(abc)^{-2}<0$,
$y_6=(abc)^{-2}>0$. Corner evaluations at the corrected
constants (exact):
$w_0(\delta{=}-0.0001)=\tfrac{1153107070889}{11233957316589}$
(binding), $w_0(\delta{=}+0.0001)
=\tfrac{1151185570889}{11233957316589}$.

\emph{Construction of the atoms.} The triple
$(a,b,c)=(\tfrac{27}{50},\tfrac{263}{200},\tfrac{103}{50})$
is extracted from the primal extremal measure of the moment
problem, then rationalized; its validity does not depend on the
extraction. In detail: (i) the discretized linear program
(maximize origin mass subject to
$m_1,\dots,m_4$ pinned, $m_5\ge M_5$, $m_6\le M_6$; grid
$dx=0.001$) has optimal measure supported on four atoms,
$\{0,\ 0.5385,\ 1.3195,\ 2.0645\}$ at the binding corner ---
complementary slackness forces the optimal dual polynomial to
vanish doubly at the three non-origin atoms, which dictates the
shape $P=[(x-a)(x-b)(x-c)]^{2}/(abc)^{2}$; (ii) the atoms are
rationalized to low-height fractions within $5\times10^{-3}$ of
the extremal support; (iii) \emph{any} positive rational triple
yields a valid certificate --- $P$ is a perfect square, so
$P\ge0$; $P(0)=1$ by the normalization; the monic numerator has
$x^{5}$-coefficient $-2(a{+}b{+}c)<0$ and $x^{6}$-coefficient
$1>0$, so $y_5<0<y_6$ automatically --- hence the LP solver, the
grid, and the extraction all leave the trust base, and only the
exact rational evaluation of $w_0$ at the chosen triple remains,
which is Lemma~\ref{l:dual} (kernel-checked). The rounding cost
of the rationalization is $6\times10^{-6}$ in the headline (LP
optimum $0.794716$ at the same corner against the certified
$0.794710$). Verification:
\texttt{certification/certify\_lp.py}; Lean:
\texttt{lean/RhGate/Certificate.lean}, kernel-compiled
(\S\ref{s:verif}(f)): the denominator-cleared identity, this
monic expansion, nonnegativity, both corners, the
corner-selection sign $y_5+6y_6<0$ and the headlines; zero
\texttt{sorry}s.

\begin{thebibliography}{99}

\bibitem{BGST} S.~A.~C.~Baluyot, D.~A.~Goldston,
A.~I.~Suriajaya, and C.~L.~Turnage-Butterbaugh, \emph{Pair
correlation of zeros of the Riemann zeta function I: proportions
of simple zeros and critical zeros}, arXiv:2501.14545 (2025).

\bibitem{GS} D.~A.~Goldston and A.~I.~Suriajaya, \emph{Zeta zeros
on the critical line}, arXiv:2511.20059 (2025).

\bibitem{BCY11} H.~M.~Bui, J.~B.~Conrey, and M.~P.~Young,
\emph{More than 41\% of the zeros of the zeta function are on the
critical line}, Acta Arith.\ \textbf{150} (2011), 35--64.

\bibitem{C26} Claude (Anthropic), \emph{More than two thirds
of the zeros of the Riemann zeta function lie on the critical
line}, preprint, August 2026,
\url{https://www-cdn.anthropic.com/564f962e60643842f5fcb4a17c9dbc8f608f1c37.pdf}.

\bibitem{Con89} J.~B.~Conrey, \emph{More than two fifths of the
zeros of the Riemann zeta function are on the critical line},
J.\ reine angew.\ Math.\ \textbf{399} (1989), 1--26.

\bibitem{Dav} H.~Davenport, \emph{Multiplicative Number Theory},
3rd ed., Springer GTM 74, 2000.

\bibitem{Fen12} S.~Feng, \emph{Zeros of the Riemann zeta function
on the critical line}, J.\ Number Theory \textbf{132} (2012),
511--542.

\bibitem{HR} H.~Halberstam and H.-E.~Richert, \emph{Sieve
Methods}, Academic Press, 1974.

\bibitem{IK} H.~Iwaniec and E.~Kowalski, \emph{Analytic Number
Theory}, AMS Colloquium Publ.\ 53, 2004.

\bibitem{Lev74} N.~Levinson, \emph{More than one third of zeros
of Riemann's zeta-function are on $\sigma=1/2$}, Adv.\ Math.\
\textbf{13} (1974), 383--436.

\bibitem{MRT} K.~Matom\"aki, M.~Radziwi\l{}\l, and T.~Tao,
\emph{Correlations of the von Mangoldt and higher divisor
functions I. Long shift ranges}, Proc.\ London Math.\ Soc.\ (3)
\textbf{118} (2019), 284--350.

\bibitem{Mon73} H.~L.~Montgomery, \emph{The pair correlation of
zeros of the zeta function}, Proc.\ Sympos.\ Pure Math.\
\textbf{24} (1973), 181--193.

\bibitem{Mon75} H.~L.~Montgomery, \emph{Distribution of the zeros
of the Riemann zeta function}, Proc.\ ICM Vancouver 1974, vol.~1,
379--381.

\bibitem{PRZZ} K.~Pratt, N.~Robles, A.~Zaharescu, and
D.~Zeindler, \emph{More than five-twelfths of the zeros of
$\zeta$ are on the critical line}, Res.\ Math.\ Sci.\ \textbf{7}
(2020), no.~2.

\bibitem{Sel42} A.~Selberg, \emph{On the zeros of Riemann's
zeta-function}, Skr.\ Norske Vid.-Akad.\ Oslo I (1942), no.~10.

\bibitem{Tit86} E.~C.~Titchmarsh, \emph{The Theory of the
Riemann Zeta-Function}, 2nd ed.\ (revised by D.~R.~Heath-Brown),
Oxford University Press, 1986.

\bibitem{Wei52} A.~Weil, \emph{Sur les ``formules explicites''
de la th\'eorie des nombres premiers}, Comm.\ S\'em.\ Math.\
Univ.\ Lund (1952), 252--265.

\end{thebibliography}
\end{document}
